% 本文件由 LaTeX 排版 Agent 根据有效 translation.json 自动生成。
% author=Theophilus; work_id=0204; title=Theophilus致Autolycus

\chapter{\Person{Theophilus}致\Person{Autolycus}}

% structure: p0001 source=h1 target=omitted reason=page-title-already-rendered-by-parent

\Emph{第一卷}\newline{}\Emph{第二卷}\newline{}\Emph{第三卷}\par

\SourceRecord{Translated by Marcus Dods.；Ante-Nicene Fathers；Vol. 2.；Edited by Alexander Roberts, James Donaldson, and A. Cleveland Coxe.；Buffalo, NY: Christian Literature Publishing Co.,；1885.；<http://www.newadvent.org/fathers/0204.htm>.}

% structure: p0001 source=h1 target=section reason=page-title

\section{\WorkTitle{致奥托吕科斯，卷一}}

% structure: p0002 source=h2 target=subsection reason=relative-heading-rank; base=h2

\subsection{第一章　奥托吕科斯——偶像崇拜者与基督徒的蔑视者}

流利的言辞与优雅的文风，为那些心智败坏的可悲之人带来愉悦和虚荣所喜爱的赞美；但喜爱真理的人并不留意华丽的辞藻，而是考察言论的实际内容，看它是什么，以及属于何种性质。因此，我的朋友，你既以空言攻击我，夸耀你那用木头和石头、锤打和铸造、雕刻和镂刻的神明，它们既不能看，也不能听，因为它们是偶像，是人手所做的作品；此外，你还称我为基督徒，仿佛这是一个该受诅咒的名号。我则声明我是基督徒，并怀着为天主服务的希望，承担这蒙天主所爱的名号。因为并不像你所想像的那样，天主的名字是难以承担的；但你之所以对天主抱有这种看法，或许是因为你自己尚未服务于祂。\par

% structure: p0004 source=h2 target=subsection reason=relative-heading-rank; base=h2

\subsection{第二章　必须先净化灵魂的眼睛，才能看见天主}

但若你问：\Quote{请把你的天主指给我看}，我要回答：\Quote{请把你自己指给我看，我便将我的天主指给你看。}那么，请显示你灵魂的眼睛能够看见，你心中的耳朵能够听见；因为正如人用肉体的眼睛看待地上的事物和今生的事，同时分辨差异——光明或黑暗、白或黑、丑或美、匀称相称或不成比例和笨拙、畸形或残缺；同样，借着听觉我们分辨尖锐、低沉或悦耳的声音；同样，对于灵魂的眼睛和心中的耳朵，我们能够借此看见天主。因为天主被那些能看见祂的人所看见，当他们开启灵魂的眼睛时：因为所有人都拥有眼睛，但在某些人那里，眼睛被遮蔽，看不见太阳的光。然而，盲人看不见并不等于太阳的光不照耀；让盲人责备自己和自己的眼睛吧。同样，人啊，你灵魂的眼睛被你的罪恶和恶行所遮蔽。人应当像一面磨光的镜子一样保持灵魂纯净。当镜子上有了锈迹，就不可能看清镜中人的面容；同样，当人心中有罪时，这样的人就不能看见天主。因此，请将你自己指给我看，看看你是否不是奸淫者、淫乱者、偷窃者、强盗、盗贼；是否不污秽少年；是否不傲慢、诽谤、暴躁、嫉妒、骄傲、目空一切；是否不争吵、贪婪、不孝敬父母；是否不卖自己的孩子。因为对于做这些事的人，天主是不显现的，除非他们先清除自己的一切不洁。所有这些事都使你陷入黑暗，正如眼睛上的薄膜阻碍人看见太阳的光；同样，不义之事，人啊，使你陷入黑暗，以致你不能看见天主。\par

% structure: p0006 source=h2 target=subsection reason=relative-heading-rank; base=h2

\subsection{第三章　天主的本性}

那么你会对我说：\Quote{你既然看见天主，请告诉我天主的容貌。}请听，人啊。天主的容貌是不可言说、不可描述的，不能以肉体的眼睛看见。因为祂在荣耀上是不可理解的，在伟大上是不可测度的，在高超上是不可思议的，在权能上是无可比拟的，在智慧上是无与伦比的，在良善上是无可模仿的，在仁慈上是无法言说的。因为若我说祂是光，我所说的不过是祂的工程；若我称祂为圣言，我所说的不过是祂的统治；若我称祂为理智，我所说的不过是祂的智慧；若我说祂是灵，我所说的是祂的气息；若我称祂为智慧，我所说的是祂的后裔；若我称祂为力量，我所说的是祂的权柄；若我称祂为权能，我是在提及祂的行动；若我称祂为眷顾，我不过是在提及祂的良善；若我称祂为王国，我不过是在提及祂的荣耀；若我称祂为主，我是在提及祂是审判者；若我称祂为审判者，我是说祂是公义的；若我称祂为父，我是说万物出于祂；若我称祂为火，我不过是在提及祂的愤怒。那么你会对我说：\Quote{天主发怒吗？}是的，祂向作恶的人发怒，但祂对爱祂和敬畏祂的人是良善、仁慈、怜悯的；因为祂是虔敬者的惩戒者，义人的父亲；但对不敬虔者，祂是审判者和惩罚者。\par

% structure: p0008 source=h2 target=subsection reason=relative-heading-rank; base=h2

\subsection{第四章　天主的属性}

祂是无始的，因为祂是非受生的；祂是不变的，因为祂是不朽的。祂被称为\OriginalTerm{天主}{\GreekText{Θεός}}，因为祂将万物放置\OriginalTerm{放置}{\GreekText{τεθεικέναι}}在自己所提供的保障上；又因\OriginalTerm{奔跑}{\GreekText{θέειν}}，因为\OriginalTerm{奔跑}{\GreekText{θέειν}}意为奔跑、运动、活跃、滋养、预见、治理并使万物活起来。但祂是主，因为祂统治宇宙；是父，因为祂在万有之先；是塑造者与制造者，因为祂是宇宙的创造者和制造者；是至高者，因为祂在万有之上；是全能的，因为祂自己统治并包容万有。因为高天、深渊的底部和地极都在祂手中，没有祂安息之处。因为诸天是祂的工程，大地是祂的造化，海洋是祂的手工；人是祂的塑造和肖像；太阳、月亮和星辰是祂的元素，为记号、季节、日子和年岁而造，为要服事并作人的奴仆；天主从无中创造了万有，使之成为存有，为要藉着祂的工程，使祂的伟大得以认识并理解。\par

% structure: p0010 source=h2 target=subsection reason=relative-heading-rank; base=h2

\subsection{第五章　藉着受造之物认识不可见的天主}

因为人的灵魂看不见，对人而言是不可见的，却通过身体的运动被感知，同样，天主确实不能被人的肉眼看见，却通过祂的眷顾和作为被看见和感知。因为，就像任何人看见一艘装备齐全、扬帆航行、驶向港口的海船时，无疑会推断船上有一位舵手在驾驶它；同样，我们必须认识到天主是整个宇宙的治理者，尽管祂对肉身的眼睛是不可见的，因为祂是不可理解的。因为如果一个人不能直视太阳，尽管它是一个非常小的天体，因其极度的热力和力量，那么一个凡人岂不更不能面对那不可言喻的天主的荣耀吗？因为正如石榴，外皮包裹着它，里面有许多被薄膜隔开的细胞和隔间，并且有许多籽粒住在其中，同样，整个受造物被天主的神所包容，而包容的神连同受造物一起被天主的手所包容。因此，正如住在里面的石榴籽粒无法看见外皮之外的东西，因为它自己在里面；同样，人连同整个受造物都被天主的手所包围，不能看见天主。再者，一位地上的君王，即使不被所有人看见，也被相信存在；因为他通过他的法律、法令、权柄、武力和雕像被认识；而你却不愿意天主通过祂的作为和大能举动被认识吗？\par

% structure: p0012 source=h2 target=subsection reason=relative-heading-rank; base=h2

\subsection{第六章　天主通过祂的作为被认识}

人哪，观察祂的作为——季节的适时更替和温度的变化；星辰的有规律运行；昼夜、月份和年份的有序交替；种子、植物和果实的多样美丽；以及四足动物、飞鸟、爬行动物和鱼类（包括河流和海中的）的多种物种；或者观察这些动物中被赋予的生养后代的本能，不是为了它们自己的利益，而是为了人的使用；以及天主为所有血肉提供食物的眷顾，或者祂所规定所有事物服务于人类的隶属关系。观察甜美的泉源和永不枯竭的河流的流淌，以及露水、阵雨和雨水的及时供应；天体的多样运动，晨星升起、预告完美发光体的临近；以及昴宿星团、猎户座、大角星，以及其他在天空中盘旋的星辰的轨道，所有这些天主的多样智慧都以它们自己的名字称呼它们。唯独祂是天主，祂从黑暗中造光，从祂的宝库中引出光，并形成了南风的宝库，\\BibleRef{约 9:9} 以及深渊的宝库、海的边界、雪和冰雹的宝库，把水收集在深渊的仓库中，把黑暗收集在祂的宝库中，并从祂的宝库中引出甜美的、可心的、愉快的亮光；\Quote{祂使云雾从地极上升；祂为雨造闪电；}祂发出雷声以惊吓，并通过闪电预告雷的轰鸣，使任何灵魂不会因突然的震动而晕厥；而祂又如此缓和从天上闪出的闪电的威力，使它不烧毁大地；因为，如果闪电被允许其全部力量，它会烧毁大地；如果雷被允许其全部力量，它会推翻其中所有的工作。\par

% structure: p0014 source=h2 target=subsection reason=relative-heading-rank; base=h2

\subsection{第七章　当我们穿上不朽时，将看见天主}

这是我的天主，万有的主，唯独祂铺张了诸天，并在其下确立了大地；祂搅动海的最深处，使它的波涛咆哮；祂统治它的力量，平静它波浪的骚动；祂在水上奠定了大地，并赐予一种精神来滋养它；祂的气息给万物带来光明，如果祂收回祂的气息，万物将完全灭亡。藉着祂，你说话，人哪；你呼吸祂的气息，却不知道祂。这就是你的状况，因为你灵魂的瞎眼和心的刚硬。但是，如果你愿意，你可以被治愈。把自己交给那位医生，祂会为你灵魂和心灵的眼睛作白内障手术。谁是那位医生？就是天主，祂通过祂的圣言和智慧治疗并赐予生命。天主藉着祂自己的圣言和智慧创造了万物；因为\Quote{藉着祂的圣言，诸天被造，诸天万象藉着祂口中的气被造。}祂的智慧最为卓越。天主以智慧建立了大地；以知识预备了诸天；以悟性使大渊的泉源裂开，云彩降下露水。如果你领悟了这些事，人哪，过着贞洁、圣洁、公义的生活，你就能看见天主。但在一切之前，让信德和敬畏天主在你心中掌权，然后你就会明白这些事。当你脱去可朽的，穿上不朽的，那时你就能相称地看见天主。因为天主将使你的肉身与你的灵魂一同复活，成为不朽的；然后，成为不朽的你将看见那不朽者，如果你现在相信祂；那时你将知道你曾不公正地说了反对祂的话。\par

% structure: p0016 source=h2 target=subsection reason=relative-heading-rank; base=h2

\subsection{第八章　一切事情都需要信德}

但你不相信死人复活。当复活发生时，那时你将要相信，不管你愿意与否；你的不信将被算作不信，除非你现在相信。你为什么不信呢？难道你不知道信德在一切事情中是首要原则吗？哪个农夫能收获，除非他先把种子信托给土地？谁能渡海，除非他先把自己信托给船和舵手？哪个病人能被治愈，除非他先把自己信托给医生的照顾？哪门艺术或知识能学会，除非他先向师傅申请并信托给他？那么，如果农夫信托了土地，水手信托了船，病人信托了医生，你难道还不信任天主，即使你手里持有这么多来自祂的保证？因为首先祂从无中创造了你，并把你带入存在（因为如果你的父亲不存在，你的母亲也不存在，那么你自己更是曾经不存在），并从微小潮湿的质料，甚至从最小的滴点中形成了你，这滴点本身一度不存在；天主把你引入此生。而且，你相信人手所造的像是神，并做大事；难道你不能相信那造你的天主也能以后再造你吗？\par

% structure: p0018 source=h2 target=subsection reason=relative-heading-rank; base=h2

\subsection{第九章 众神的邪淫。}

而且，确实，你声称敬拜的那些人的名字，是死人的名字。而这些人是怎样的人、什么样的人呢？萨图尔努斯不是被发现是食人者，毁灭并吞吃自己的儿女吗？如果你称他的儿子为尤庇特，也听听他的事迹和行为——首先，他如何在伊达山上被一只母山羊哺乳，并根据神话，杀死了它，剥了它的皮，用兽皮做了一件外套。他的其他事迹——他的乱伦、奸淫和贪欲——荷马和其他诗人将会更好地叙述。我何必再提他的儿子们？赫丘利如何自焚；以及醉醺醺狂乱的巴克斯；阿波罗害怕并逃离阿基琉斯，爱上达芙妮，不知道雅辛托斯的命运；维纳斯受伤，玛尔斯是凡人的灾祸；还有从所谓的神身上流出的\OriginalTerm{神血}{ichor}。而这些确实是传说中较温和的种类；因为被称为奥西里斯的神被发现肢体被撕裂，每年庆祝他的奥秘，好像他死了，又被找到，一块块地寻找。因为既不知道他是否死了，也没有显示他是否被找到。我何必再提被阉割的阿提斯，或在林中游荡、打猎时被野猪刺伤的阿多尼斯；或被雷电击中的埃斯库拉庇乌斯；或从锡诺普逃到亚历山大的逃亡者塞拉比斯；或西徐亚的狄安娜，她自己也是逃亡者、杀人和女猎手，以及对恩底弥翁的热恋？现在，不是我们宣扬这些事，而是你们自己的作家和诗人。\par

% structure: p0020 source=h2 target=subsection reason=relative-heading-rank; base=h2

\subsection{第十章 偶像崇拜的荒谬。}

我何必再详述埃及人所崇拜的大量动物，包括爬行动物、牛、野兽、鸟类、河鱼；甚至洗盆和可耻的声响？但如果你援引希腊人和其他民族，他们崇拜石头、木头和其他种类的物质材料——正如我们刚才所说的，死人的像。因为斐迪阿斯被发现，在皮萨为埃利斯人制造了奥林匹亚的尤庇特，在雅典制造了雅典卫城的密涅瓦。我要问你，我的朋友，存在多少尤庇特？因为首先有被称为奥林匹亚的尤庇特，然后有拉提亚里斯尤庇特、卡西乌斯尤庇特、雷神尤庇特、先祖尤庇特、全夜尤庇特、城邦尤庇特、卡皮托利尼尤庇特；还有那个克里特人国王、萨图尔努斯之子的尤庇特，在克里特有坟墓，但其余的，可能被认为不配得到坟墓。如果你提到那被称为众神之母的，我绝不敢用我的嘴唇说出她的事迹，或那些敬拜她的人的事迹（因为对我们来说，甚至提到这些事都是不合法的），以及她和她的儿子们向国王提供多少税收和收入。因为这些不是神，而是偶像，正如我们已经说过的，是人手的工作和不洁的魔鬼。但愿所有制造它们并信赖它们的人也变成这样！\par

% structure: p0022 source=h2 target=subsection reason=relative-heading-rank; base=h2

\subsection{第十一章 国王应受尊敬，天主应受敬拜。}

因此，我宁愿尊敬国王[胜过你们的神]，但不是敬拜他，而是为他祈祷。但我敬拜生活的、真实的天主，知道国王是由祂设立的。那么，你会对我说：\Quote{你为什么不敬拜国王？}因为他被造不是为了被敬拜，而是应当以合法的尊敬来敬畏，因为他不是神，而是天主所任命的人，不是为被敬拜，而是为公正地审判。因为在某种程度上，他的治理是由天主委托给他的：正如祂不会让那些祂在自己以下所任命的人被称为国王；因为\Quote{国王}是他的称号，别人使用它是不合法的；同样，除了天主以外，任何人被敬拜也是不合法的。因此，人啊，你完全错了。因此，要尊敬国王，顺服他，以忠诚的心为他祈祷；因为如果你这样做，你就是遵行天主的旨意。因为天主的法律说：\Quote{我儿，敬畏主和国王，不要违抗他们；因为他们要突然向他们的仇敌施行报复。}\par

% structure: p0024 source=h2 target=subsection reason=relative-heading-rank; base=h2

\subsection{第十二章 基督徒名称的含义。}

关于你嘲笑我并称我为\Quote{\OriginalTerm{基督徒}{Christian}}，你不知道你在说什么。首先，因为那受傅的是甘甜和有用的，远非可鄙视的。因为哪艘船能有用且经得起海，除非它先被填缝（受傅）？哪座城堡或房屋，在没有受傅时是美丽和有用的？哪个人，当他进入此生或进入运动场时，不是用油受傅的？什么作品有装饰或美丽，除非它被受傅和抛光？其次，空气和天下万物在某种程度上被光和灵受傅；而你难道不愿意用天主的油受傅吗？因此，因着这个原因，我们被称为\OriginalTerm{基督徒}{Christian}，因为我们用天主的油受傅。\par

% structure: p0026 source=h2 target=subsection reason=relative-heading-rank; base=h2

\subsection{第十三章 由例证证明的复活。}

那么，关于你否认死人\OriginalTerm{复活}{resurrection}的事——因为你说：\Quote{请向我展示一个从死者中复活的人，让我亲眼看见而相信}——首先，你若看见了事情的发生才相信，这有什么了不起呢？再者，你相信自焚的\Person{Hercules}仍然活着；你相信被闪电击中的\Person{Æsculapius}复活了；而对于天主告诉你的事，你却不相信吗？但是，假设我向你展示一个死人复活而活着，即使这样你也不会相信。天主确实向你展示许多凭据，好让你相信祂。因为，请想一想，季节、白天和黑夜的消逝，它们如何死亡又重新升起。还有，难道不是有种子和果实的复活，且是为人的使用吗？例如，一颗麦子或其他谷物的种子，被撒在地里，先死去腐烂，然后复活，成为麦秆。树木和果树的本质，不正是按照天主的安排，从不可见、不可视中按其季节结出果实吗？此外，有时一只麻雀或其他鸟类，喝水时吞下了一颗苹果或无花果的种子，或其他东西，飞到某座岩石小山或坟墓上，将种子留在粪便中；这颗曾被吞下、经过如此高温的种子，如今扎根，长成了一棵树。所有这些都是天主的智慧所成就的，为要藉着这些事显明，天主能够实现全人类的普遍复活。若你愿见证一个更奇妙的现象，这现象不仅证明地上之物的复活，也证明天上之物的复活，请思考月亮的复活，它每月发生：如何亏缺、死亡，又重新升起。人啊，请进一步听你自己身上进行的复活工程，尽管你未曾察觉。因为你或许曾患病，失去了肉体、力量和美丽；但当你从天主那里再次获得怜悯和医治时，你又重新获得了肉体和容貌，并恢复了力量。正如你不知道你的肉体去往何处消失，你也不知道它从哪里生长、从哪里再来。但你会说：\Quote{从饮食变化成血。}确实如此；但这也是天主的工作，祂如此运作，而非任何其他。\par

% structure: p0028 source=h2 target=subsection reason=relative-heading-rank; base=h2

\subsection{第十四章 \Person{特奥菲卢斯}——归化的榜样}

所以，不要怀疑，而要相信；因为我自己也曾不信这事会发生，但现在，经过思考这些事，我相信了。同时，我遇见了神圣先知们的圣经，他们也是藉着天主的圣神预言了已经发生的事，正如它们所实现的，以及正在发生的事，正如它们正在发生的，还有未来的事，按它们将要成就的顺序。因此，既然承认事件按预言发生所提供的证据，我不不信，反而相信，顺从于天主；你若愿意，也请服从祂，相信祂，免得如今你若继续不信，日后当你在永罚中受折磨时才被说服；这些刑罚，当它们被先知预言后，后来的诗人和哲学家从圣经中窃取，以使他们的教义值得信任。然而，他们也预先讲述了将要降临在亵渎者和不信者身上的刑罚，为使人没有留下的见证，或能说：\Quote{我们没有听说过，也不知道。}但你也若愿意，请恭敬地留意先知书，它们将使你的道路更平坦，以逃避永罚，并获得天主的永恒奖赏。因为那位赐口舌以说话、造耳朵以听、造眼睛以看的天主，将审查一切，并施行公义的审判，给每人应得的报应。对于那些藉着恒心行善\BibleRef{罗 2:7}寻求不朽的人，祂将赐予永生、喜乐、平安、安息，以及丰盛的美物，是眼睛未曾看见，耳朵未曾听见，人心也未曾想到的。\BibleRef{格前 2:9}但对于不信者和藐视者，他们不顺从真理，反而顺从于不义，当他们充满了奸淫、邪淫、污秽、贪婪和不合法的偶像崇拜时，将有愤怒和忿怒、患难和困苦，\BibleRef{罗 2:8}最后，永火将占据这样的人。既然你说：\Quote{把你的天主指给我看，}这就是我的天主，我劝你敬畏祂，信赖祂。\par

\SourceRecord{Translated by Marcus Dods.；Ante-Nicene Fathers；Vol. 2.；Edited by Alexander Roberts, James Donaldson, and A. Cleveland Coxe.；Buffalo, NY: Christian Literature Publishing Co.,；1885.；<http://www.newadvent.org/fathers/02041.htm>.}

% structure: p0001 source=h1 target=section reason=page-title

\section{致奥托里库斯，第二卷}

% structure: p0002 source=h2 target=subsection reason=relative-heading-rank; base=h2

\subsection{第一章 撰写本书的缘由}

从前，我很好的朋友奥托里库斯，我们曾有过一些交谈，当你问我的天主是谁，并稍加留意我的言论时，我向你解释了我的宗教信仰；然后彼此告别，我们带着许多相互的友好之情各自回家，尽管起初你对我有些不客气。因为你知道且记得，你曾以为我们的教义是愚拙。后来你催促我这样做，我虽然未曾受过演讲术的训练，但渴望通过这篇论文更准确地证明你所陷于的虚妄劳作和空洞崇拜；并且我也希望，从你阅读、或许尚未完全理解的你们自己的少数历史中，向你显明真理。\par

% structure: p0004 source=h2 target=subsection reason=relative-heading-rank; base=h2

\subsection{第二章 神明被造时受轻视，被购买时却变得珍贵}

真的，在我看来，雕塑家、雕刻家、画家或塑像者，竟然能够设计、绘画、雕刻、塑造和制作神明，这些神明被工匠造出时被认为毫无价值；但一旦被某人购得并放置在某个所谓的庙宇或某所房屋中，不仅那些购买者向他们祭献，连那些制造和出售他们的人也带着极大的虔诚、祭献的器具和奠祭前来敬拜；他们视之为神明，却未看出它们与当初被自己制造时一模一样，不论材料是石头、铜、木头、颜料或其他什么东西。你的情况也是如此，当你阅读所谓神明的历史和谱系时。因为你读到他们的出生，便以为他们是人，但后来又称他们为神明并敬拜他们，没有反思也不明白，他们生时正是你所读到的那样。\par

% structure: p0006 source=h2 target=subsection reason=relative-heading-rank; base=h2

\subsection{第三章 神明究竟怎么了？}

至于古代的神明，如果他们确实曾被生养，那么繁衍是相当多的。但现在，他们的生育在哪里显明呢？因为如果古时他们生育并受生，显然直到如今也应该有神明被生养出生；至少如果不是这样，这样的种族就会被视为无能。因为他们要么已经衰老，因此不再生育，要么已经灭绝，不复存在。因为如果神明是被生的，他们理当直到现在也仍在出生，正如人类也被出生一样；是的，神明应该比人类更多，正如西彼尔所说：——\par

\Quote{因为神明若生育，\newline{}且各各永存，\newline{}那么神明的种族必多于凡人，\newline{}其数量如此之多，以至凡人没有站立之处。}\par

因为如果那些必死短命之人的儿女直到如今还出现，人类并没有停止出生，以致城市和村庄都满了，甚至乡野也有人居住，那么那些据你们诗人所说不会死亡的神明，岂不更应该生育和受生吗？既然你们说神明是由生育产生的。还有，那被称为奥林帕斯的山从前曾由神明居住，现在为何荒废？或者，为何古时的朱庇特曾住在伊达山，并据荷马和其他诗人所说，被人知道住在那里，但现在却无人知晓？为何他只在世界的一个部分被找到，而不是无处不在？因为他要么忽略了其他部分，要么不能无所不在并眷顾一切。例如，如果他在东方，他就不能在西方；反之，如果他在西方，他就不在东方。但这乃是至高全能者、永生天主的属性：不仅无所不在，而且看见一切、听见一切，丝毫不受空间限制；因为如果祂被空间限制，那么包含祂的空间就比祂更大；因为包含者大于被包含者。因为天主不被包含，反而祂是万有的空间。但朱庇特为何离开了伊达山？是因为他死了，还是那座山不再使他喜悦？他去了哪里？到天上？不。但或许你会说，到克里特岛？是的，因为那里至今还展示着他的坟墓。你又会说，到皮萨，在那里直到今日他还反射着菲迪亚斯双手的光荣。那么，让我们继续看哲学家和诗人的著作。\par

% structure: p0010 source=h2 target=subsection reason=relative-heading-rank; base=h2

\subsection{第四章 哲学家们关于天主的荒谬见解}

廊下派的某些哲学家说根本没有天主；即使有，他们也说祂只关心自己而不顾其他。伊壁鸠鲁和克吕西普的愚昧已详尽阐述了这些观点。另一些人说万物无需外因产生，世界非受造，本性永恒；他们竟敢断言根本没有天主的眷顾，而主张天主不过是各人的良心。还有一些人主张遍及万物的精神就是天主。但柏拉图及其学派确实承认天主非受造，万有的圣父和创造者；但他们又主张物质与天主同为非受造，并断言物质与天主一样永恒。然而，如果天主非受造，物质也非受造，那么按照柏拉图主义者的说法，天主就不再是万物的创造者，就他们的观点而言，天主的唯一统治也无法确立。再者，正如天主因其非受造也是不变的，那么如果物质也是非受造的，它也会是不变的，并且与天主平等；因为受造之物是可变的、可更改的，而非受造之物则是不变的、不可更改的。如果天主用已有的材料创造世界，那有什么了不起呢？因为即使一个人类的工匠，从某人那里得到材料后，也能随心所欲地制造。但天主的能力正表现在这一点上：祂从无中创造出祂所喜悦的一切；正如赋予生命和行动的能力唯属于天主，而非任何其他。因为人确实能制作一个形象，却不能给所造之物理性、气息或感觉。但天主在这一点上远超人类的能力：祂能造出具有理性、生命和感觉的作品。因此，正如在所有这些方面天主比人更强有力，同样，祂从无中创造并创造了现有的万物，且随心所愿、随己意行。\par

% structure: p0012 source=h2 target=subsection reason=relative-heading-rank; base=h2

\subsection{第五章 荷马与赫西俄德关于诸神的意见。}

因此，你们那些哲学家和作家的意见是不一致的；因为前者提出了上述观点，而诗人荷马则基于完全不同的假设阐述了世界甚至诸神的起源。他在某处说道：——\par

\Quote{诸神之父俄刻阿诺斯，\newline{}以及那生育诸神的母亲忒提斯，\newline{}一切河流与每一片海洋都源自他们。}\par

然而，他这样说时，并没有把天主启示给我们。因为谁不知道大洋是水呢？但如果是水，就不是天主。事实上，如果天主是万物的创造者，祂确实是，那么祂也是水和海洋的创造者。赫西俄德本人也宣告了不仅是诸神，连世界本身的起源。尽管他说世界是被造的，却没有告诉我们是谁创造了它。此外，他说克洛诺斯及其子朱庇特、尼普顿和普路托是神，而我们发现他们是在世界之后出生的。他还讲述了克洛诺斯如何被自己的儿子朱庇特在战争中攻击；因为他说道：——\par

\Quote{他以武力战胜了其父克洛诺斯，\newline{}在永生者中公正而智慧地统治，\newline{}并给每一位分配了适当的尊荣。}\par

然后他在诗中引入了朱庇特的女儿们，他称之为缪斯，并作为她们的恳求者出现，渴望从她们那里查明万物是如何形成的；因为他说道：——\par

\Quote{朱庇特的女儿们，万福！请赐我助佑，\newline{}使我以优美有序的诗句，\newline{}歌唱不朽诸神的诞生；\newline{}他们生于繁星的天穹和大地；\newline{}他们最初从幽暗的夜中生出，\newline{}由咸海养育长大。\newline{}也请告诉我，谁首先给大地以形态，\newline{}以及河流和无边的大海，\newline{}其浪涛不倦地下沉，又高高扬起；\newline{}还有那缀满星辰的辉煌穹苍\newline{}是如何铺展开来的。\newline{}赐予一切美善的诸神又源自何处？\newline{}请说，他们手中赐下了哪些不同的礼物，\newline{}使有些人得财富，有些人得荣誉或名声；\newline{}他们从最遥远的时代起，\newline{}如何居住在阳光明媚的奥林波斯山众多洞穴中。\newline{}这些事情，你们缪斯啊，请述说，\newline{}你们永远居住在奥林波斯的阴影中——因为你们能讲述：\newline{}从起初你们的脚步就已在那里徜徉；\newline{}那就告诉我们，万物中哪一样最先被造。}\par

但缪斯们比世界年轻，又怎能知道这些事呢？当她们的父亲尚未出生时，她们又如何能向赫西俄德讲述[正在发生的事]呢？\par

% structure: p0020 source=h2 target=subsection reason=relative-heading-rank; base=h2

\subsection{第六章 赫西俄德论世界的起源。}

他在一定程度上承认物质[是自存的]，并承认世界[没有造物主]被造，说道：——\par

\Quote{首先，从混沌中造出了万物，\newline{}然后，宽胸大地的根基被牢固安置，\newline{}永生者安全地居住其上，\newline{}他们长居白雪皑皑的奥林波斯。\newline{}后来，黑暗的塔塔罗斯诞生了，\newline{}在宽阔大地的深处，\newline{}还有爱神，即使在众神中也是最美丽的，\newline{}祂征服一切，降服心灵和意志，\newline{}解除忧虑，并在神与人心中赐予\newline{}明智的忠告。\newline{}从混沌生出厄瑞玻斯和黑夜，\newline{}从黑夜和厄瑞玻斯生出大气和黎明。\newline{}大地按自己的形象造出了繁星的天穹，\newline{}为万物提供了庇护，\newline{}并使有福的神明得以安息。\newline{}凭她的力量，崇山峻岭升起，\newline{}为树神造出了宜人的洞穴，\newline{}生出了荒芜的大海及其所有波涛，\newline{}未经爱情；但当大地受到天穹的爱慕，\newline{}她便生出了深邃的漩涡大洋。}\par

他说了这番话，却尚未说明这一切是由谁造成的。因为若起初就有混沌存在，某种物质是非受造的、预先存在，那么是谁改变了它的状态，赋予它不同的秩序和形状呢？是物质自身改变了自己的形式，将自己安排成一个世界吗？（因为\Person{朱庇特}的诞生不仅远在物质之后，而且远在世界和许多人之后；他的父亲\Person{萨图尔努斯}也是如此。）还是有一位统治万物的权能者创造了它？我当然是指天主，祂也将它塑造成一个世界。此外，他在各方面都被发现胡言乱语，自相矛盾。因为当他提到地、天和海时，他让我们明白，从这些中产生了诸神；又从这些（诸神）中，他宣称产生了一些非常可怕的人——\OriginalTerm{提坦}{Titans}族和\OriginalTerm{独眼巨人}{Cyclopes}的族类，以及一大群巨人和埃及诸神——或者不如说，是虚妄之人，正如\Person{Apollonides}（别号\Person{Horapius}）在名为\WorkTitle{Semenouthi}的书中，以及在他关于埃及人的崇拜及其国王、以及他们所从事的虚妄劳作的其它历史中所提及的那样。\par

% structure: p0024 source=h2 target=subsection reason=relative-heading-rank; base=h2

\subsection{第七章 外邦人虚构的族谱}

我又何必重述那些希腊神话——关于黑暗之王\Person{普路托}，关于\Person{涅普顿}潜入海底，拥抱\Person{墨拉尼珀}并生下一个食人子的故事——或是你们的作家编入悲剧中的、关于\Person{朱庇特}众子的许多故事，他们之所以记载这些人的谱系，是因为他们是人而非神所生？而喜剧诗人\Person{阿里斯托芬}，在名为\WorkTitle{鸟}的戏剧中，大胆处理创世主题，说起初世界是从一个蛋中产生的，他说：\par

\Quote{一个由黑翼之夜产下的风蛋\newline{}起初。}\par

但是\Person{萨提鲁斯}，他在叙述亚历山大里亚家族的历史时，从\Person{菲洛帕托尔}（又名\Person{托勒密}）开始，宣称\Person{巴克斯}是他的祖先；因此\Person{托勒密}也是这个家族的创始人。\Person{萨提鲁斯}接着这样说：\Person{得伊阿尼拉}是\Person{巴克斯}和\Person{忒斯提乌斯}之女\Person{阿尔泰亚}所生；从她与\Person{朱庇特}之子\Person{赫拉克勒斯}，据我所推断，产生了\Person{许罗斯}；从他产生了\Person{克莱奥德摩斯}，从他又产生了\Person{阿里斯托马科斯}，从他又产生了\Person{特墨诺斯}，从他又产生了\Person{刻伊苏斯}，从他又产生了\Person{马戎}，从他又产生了\Person{忒斯特鲁斯}，从他又产生了\Person{阿库斯}，从他又产生了\Person{阿里斯托米达斯}，从他又产生了\Person{卡拉努斯}，从他又产生了\Person{科伊诺斯}，从他又产生了\Person{提瑞玛斯}，从他又产生了\Person{佩尔狄卡斯}，从他又产生了\Person{腓力}，从他又产生了\Person{埃罗普斯}，从他又产生了\Person{阿尔克塔斯}，从他又产生了\Person{阿明塔斯}，从他又产生了\Person{博克鲁斯}，从他又产生了\Person{墨勒阿革耳}，从他又产生了\Person{阿尔西诺厄}；从她与\Person{拉古斯}产生了\Person{托勒密·索特尔}，从他与\Person{阿尔西诺厄}产生了\Person{托勒密·欧厄尔革忒斯}，从他与\Person{贝勒尼基}（\Person{库瑞涅}王\Person{马伽}之女）产生了\Person{托勒密·菲洛帕托尔}。这样，亚历山大里亚诸王与\Person{巴克斯}的关系便如此确定了。因此，在\Person{狄俄倪索斯}部落中，存在着不同的家族：来自\Person{阿尔泰亚}的\OriginalTerm{阿尔泰亚族}{Althean}，\Person{阿尔泰亚}是\Person{狄俄倪索斯}的妻子、\Person{忒斯提乌斯}之女；还有\Person{得伊阿尼拉}家族，来自她是\Person{狄俄倪索斯}和\Person{阿尔泰亚}的女儿、\Person{赫拉克勒斯}之妻；——由此，家族也得名：\Person{阿里阿德涅}家族，来自\Person{米诺斯}之女、\Person{狄俄倪索斯}之妻\Person{阿里阿德涅}，一位恭敬的女儿，她以另一种形态与\Person{狄俄倪索斯}交往；\OriginalTerm{忒斯提亚族}{Thestian}，来自\Person{阿尔泰亚}之父\Person{忒斯提乌斯}；\OriginalTerm{托阿斯族}{Thoantian}，来自\Person{狄俄倪索斯}之子\Person{托阿斯}；\OriginalTerm{斯塔菲洛斯族}{Staphylian}，来自\Person{狄俄倪索斯}之子\Person{斯塔菲洛斯}；\OriginalTerm{欧诺斯族}{Euænian}，来自\Person{狄俄倪索斯}之子\Person{欧诺斯}；\OriginalTerm{马戎族}{Maronian}，来自\Person{阿里阿德涅}和\Person{狄俄倪索斯}之子\Person{马戎}；——因为所有这些人都属于\Person{狄俄倪索斯}之子。确实，许多其他名字由此起源，并流传至今，如\Person{赫拉克勒斯}族的\OriginalTerm{赫拉克利代}{Heraclidæ}，\Person{阿波罗}族的\OriginalTerm{阿波罗尼代}{Apollonidæ}，\Person{波塞冬}族的\OriginalTerm{波塞多尼}{Poseidonii}，以及源自\Person{宙斯}的\OriginalTerm{狄}{Dii}和\OriginalTerm{狄奥格纳}{Diogenæ}。\par

% structure: p0028 source=h2 target=subsection reason=relative-heading-rank; base=h2

\subsection{第八章 论天主的眷顾}

我又何必继续列举如此众多的名字和谱系呢？因此，所有的作者、诗人以及所谓哲学家，都完全被欺骗了；听从他们的人也是如此。因为他们大量编造关于他们诸神的神话和愚蠢故事，没有把他们展示为神，而是展示为人，而且是某些醉酒、淫乱和杀人之辈。但关于世界的起源，他们也发表了矛盾而荒谬的意见。首先，有些人，正如我们之前解释的，主张世界是非受造的。而那些说世界是非受造和自生的人，又反驳了那些主张世界是被造的人。因为他们凭推测和人的概念说话，并不认识真理。还有另外一些人，说存在着\OriginalTerm{天主的眷顾}{providence}，从而推翻了前人的立场。\Person{阿拉托斯}确实说：\par

\Quote{从\Person{朱庇特}开始我的歌；愿他的名字\newline{}永不缄默：一切充满你；\newline{}人的道路和居所；天和海：\newline{}我们的存在系于你；我们内行动；\newline{}全是你的后裔和\Person{朱庇特}的种子。\newline{}仁慈的祂规劝人类向善，\newline{}激励劳作，促发食物的希望。\newline{}祂指示牛群最佳放牧处，\newline{}深耕的土壤将结出黄谷。\newline{}农夫应当种植或播种何时，\newline{}由祂来告知，唯祂知道。}\par

那么，我们该相信谁呢：是此处所引的\Person{阿拉托斯}，还是\Person{索福克勒斯}，当他说：——\par

\Quote{对于未来没有预见；\newline{}最好随意地生活，随遇而安}？\par

而\Person{荷马}又不赞同这一点，因为他说美德\par

\Quote{在人类中随\Person{朱庇特}的旨意而增减。}\par

而\Person{西蒙尼德斯}说：——\par

\Quote{无人也无城邦有美德，除非来自天主；\newline{}智谋居于天主内；可怜的人\newline{}除了自己的可怜外一无所有。}\par

同样，\Person{欧里庇得斯}：——\par

\Quote{离了天主，人一无所是。}\par

而\Person{米南德}：——\par

\Quote{除天主外，无人为我们提供。}\par

\Person{欧里庇得斯}又说：——\par

\Quote{因为当天主意愿拯救时，祂将使万物弯曲\newline{}作为工具来成就祂的目的。}\par

而\Person{忒斯提乌斯}：——\par

\Quote{若天主意旨拯救你，你必安全，\newline{}即使在大洋中乘席航行。}\par

他们说了无数类似的话，却又自相矛盾。至少\Person{索福克勒斯}，他在别处否认了天主的眷顾，却说：——\par

\Quote{无人能逃脱天主的打击。}\par

此外，他们既引入了众多神灵，却又谈论独一者；他们反对那些肯定天主的眷顾的人，坚持相反的意见说没有天主的眷顾。因此\Person{欧里庇得斯}说：——\par

\Quote{我们劳苦费力，枉然耗尽力量，\newline{}因为空虚的盼望，而非远见，是我们的向导。}\par

他们无意中承认自己不认识真理；但受魔鬼的感动并被其鼓胀，便按其唆使说出了所说的一切。实际上，诗人——荷马和赫西俄德——据说是受缪斯默感，从欺骗性的幻想发言，而非出于纯洁的灵，而是出于谬误的灵。这一点从以下事实明显可见：直到今日，附魔者有时仍藉着永生的真天主之名被驱魔；而这些谬误之灵自己也承认，他们是魔鬼，也曾启发过这些作者。但有时他们中间有些人在灵魂中醒悟，为使自己成为对自身及万民的见证，便说了与先知们和谐一致的话，关于天主的独尊、审判及类似之事。\par

% structure: p0050 source=h2 target=subsection reason=relative-heading-rank; base=h2

\subsection{第九章 圣神所默感的先知}

但属于天主的人们，内携圣神而成先知，受天主默感并获得智慧，成为受天主教导的、圣洁而正义的人。因此，他们也堪当获得这一赏报：成为天主的器皿，保有来自祂的智慧，藉此智慧他们论及了世界的创造及一切其他事物。因为他们也预言瘟疫、饥荒和战争。并且不是一两位，而是许多位，在不同时期与季节中出现在希伯来人中间；在希腊人中也有西彼尔；他们所说的一切彼此一致且和谐，包括他们之前发生的事、他们自己时代发生的事，以及在我们今天正在应验的事；因此我们也确信关于未来之事将会应验，正如先前的事已经应验一样。\par

% structure: p0052 source=h2 target=subsection reason=relative-heading-rank; base=h2

\subsection{第十章 天主藉圣言创造世界}

首先，他们一致教导我们：天主从无中创造了万有；因为没有什么是与天主同在的：祂就是自己的处所，一无所缺，在万世之前就已存在，愿意造人，好让人认识祂；为此，祂为人准备了世界。因为受造者有所需求，而非受造者则一无所缺。天主在其自己的内怀中拥有祂的圣言，便生出了祂，在万有之前将祂连同祂的智慧一同发出。祂以这圣言作为祂所造万物的助手，并藉着祂创造了万有。祂被称为\Quote{\OriginalTerm{治理之本原}{\GreekText{ἁρκή}}}。祂既是天主的灵、治理之本原、智慧以及至高者的能力，便降临在先知们身上，藉着他们讲述世界的创造及其他一切。因为世界出现时，先知尚未存在，但天主的智慧在祂内，祂的圣言也常与祂同在。因此祂藉着先知撒罗满这样说：\Quote{当祂预备诸天时，我在那里；当祂奠定大地根基时，我与他同在，如同被养育的。}至于梅瑟，他生活在撒罗满之前许多年，更确切地说，是天主的话藉着他如同藉着一件工具，说道：\Quote{起初天主创造了天地。}他先提到了\Quote{起初}和\Quote{创造}，然后才介绍天主；因为不可轻率或随意称呼天主。神圣的智慧预知有些人会胡言乱语，称呼许多不存在的神。因此，为使永生的天主能藉祂的作为被认识，并使人知道天主藉祂的圣言创造了诸天和大地及其中的一切，他说：\Quote{起初天主创造了诸天和大地。}然后，在论及它们的创造之后，他向我们解释：\Quote{大地还是混沌空虚，深渊上还是一团黑暗，天主的神在水面上运行。}这正是圣经一开始所教导的，为表明天主用以制造并塑造世界的物质，在某种程度上也是受造的，是由天主产生的。\par

% structure: p0054 source=h2 target=subsection reason=relative-heading-rank; base=h2

\subsection{第十一章 记述六日工程}

现在，创造的开端是光；因为光显现了被造之物。因此经上说：\BibleQuote{天主说：\InnerQuote{要有光！}就有了光；天主看见光好，}显然是为人类而造的。\BibleQuote{天主将光与黑暗分开，称光为\InnerQuote{昼}，称黑暗为\InnerQuote{夜}。过了晚上，过了早晨，这是第一天。天主说：\InnerQuote{在水与水之间要有穹苍，将水上下分开！}事就这样成了。天主造了穹苍，把穹苍以下的水和穹苍以上的水分开。天主称穹苍为\InnerQuote{天}。天主看了，认为好。过了晚上，过了早晨，这是第二天。天主说：\InnerQuote{天下的水应聚在一处，使旱地出现！}事就这样成了。水聚在一处，旱地出现了。天主称旱地为\InnerQuote{地}，称水汇合处为\InnerQuote{海}。天主看了，认为好。天主说：\InnerQuote{地要生出青草、结种子的蔬菜和结果子的果树，各按其类，果子的种子在内。}事就这样成了。地生出了青草、结种子的蔬菜和各种结果子的树，各按其类。天主看了，认为好。过了晚上，过了早晨，这是第三天。天主说：\InnerQuote{在天空要有光体，以分别昼夜，作为时节、日子和年岁的记号；并要在天空发光，照耀大地！}事就这样成了。天主造了两个大光体：较大的管理白天，较小的管理黑夜，并造了星辰。天主将星宿安置在天空，照耀大地，管理昼夜，分别明暗。天主看了，认为好。过了晚上，过了早晨，这是第四天。天主说：\InnerQuote{水中要孳生蠕动的生物，要有飞鸟在天空飞翔！}事就这样成了。天主造了大鱼和各种水中孳生的蠕行动物，各按其类，以及各种飞鸟，各按其类。天主看了，认为好。天主祝福它们说：\InnerQuote{你们要滋生繁殖，充满海洋；飞鸟也要在地上繁殖！}过了晚上，过了早晨，这是第五天。天主说：\InnerQuote{地要生出各种生物：牲畜、爬虫和野兽！}事就这样成了。天主造了各种野兽、牲畜和爬虫。天主说：\InnerQuote{让我们照我们的肖像造人，管理海中的鱼、天空的鸟、牲畜、各种野兽和爬虫！}天主于是造了人：按天主的肖像造了人，造了一男一女。天主祝福他们说：\InnerQuote{你们要生育繁殖，充满大地，治理大地，管理海中的鱼、天空的鸟和地上爬行的生物！}天主说：\InnerQuote{看，我把地上结种子的各种蔬菜和结果子的各种树木，都给你们作食物；至于地上的野兽、天空的鸟和爬虫，我把青草作食物。}事就这样成了。天主看了他所造的一切，认为非常好。过了晚上，过了早晨，这是第六天。天地和其中的万物都完成了。在第六天，天主完成了祂所造的一切工程，就在第七天停止了一切工程，安息了。天主祝福了第七天，定为圣日，因为在这一天，祂停止了他开始创造的一切工程。}\par

% structure: p0056 source=h2 target=subsection reason=relative-heading-rank; base=h2

\subsection{第十二章 六日工程的荣耀}

对这六天的工作，没有人能对其中所有部分给出恰当的解说和描述，即使他有一万条舌头和一万张嘴；即使他在此生逗留一万年，也不能说出任何配得上这些事的话，因为上述六天工作中所含的天主智慧的极其伟大和丰富。的确，许多作家模仿了这段叙述，并试图对这些事给出解释；然而，尽管他们从中得到了一些关于创造世界和人性的暗示，却未迸发出丝毫真理的火花。哲学家、作家和诗人的言论，因其辞藻华丽而看似可信，但他们的论述被证明是愚妄而空洞的，因为其无稽轻浮之言极多，其中找不到零星半点真理。即使他们似乎说出了某些真理，也混杂着错误。正如一种有毒的药物，与蜜或酒或其他东西混合后，使整个混合物有害而无益；同样，他们口中的雄辩被发现是徒劳的，甚至对相信它的人是有害的。此外，关于第七日，人人都承认；但大多数人不知道，在希伯来人所称的\Quote{安息日}，译成希腊文就是\Quote{第七}（\OriginalTerm{第七}{\GreekText{ἑβδομάς}}），这个名称被每个民族采用，虽然他们不知道这个称呼的理由。至于诗人赫西奥德所说的从混沌中生出厄瑞玻斯，以及大地和统治其神与人的爱情，他的说法被证明是空洞、冰冷，并且完全远离真理。因为天主被快乐征服是不合适的，连有节操的人都戒除一切卑鄙的快乐和邪恶的欲望。\par

% structure: p0058 source=h2 target=subsection reason=relative-heading-rank; base=h2

\subsection{第十三章 论创造世界}

此外，他（赫西奥德）的关于天主的观念——人性的、卑下的、非常软弱的——在他开始从地上较低的事物叙述万物的创造时就暴露出来。因为人处于下方，从地开始建造，不先立根基就不能依次建成屋顶。但天主的权能正显示在此：祂首先从无中按祂的旨意创造了受造之物。\BibleQuote{因为在人不可能的，在天主却可能。}（\BibleRef{路 18:27}）因此，先知首先提到天被创造，作为一种屋顶，说道：\Quote{起初天主创造了天}——也就是说，藉着\Quote{起初}这个原理，天被造，正如我们已经指出的。而\Quote{地}他指的是土壤和根基，正如\Quote{深渊}他指的是众多的水；\Quote{黑暗}他是指天主造的天像盖子一样盖住水和大地。而运行在水面上的圣神，他指的是天主为赋予受造物生命而赐下的，如同祂将生命赐给人，将细微的与细微的混合。因为圣神是细微的，水也是细微的，这样圣神可以滋养水，而水与圣神一同渗透各处，滋养受造之物。因为圣神是唯一的，并占据光的位置，祂位于水与天之间，使得黑暗在天主说\Quote{要有光}之前，无论如何不能与更靠近天主的天空相通。因此，天像一个圆顶盖，覆盖着像土块一样的物质。于是另一位先知，名叫依撒意亚，说了这些话：\BibleQuote{是天主铺张穹苍如帐幔，展开它如可居住的帐篷。}（\BibleRef{依 40:22}）然后，天主的命令，即祂的圣言，像一盏灯在封闭的房间里照耀，照亮了天下的一切，那时祂造了光与世界分开。天主称光为昼，称黑暗为夜。因为人不可能称光为昼，或黑暗为夜，也不能给其他事物起名，除非他从创造这些事物本身的天主那里接受了命名。因此，在世界的历史和起源的最初，圣经并非指着我们所见的这个穹苍说话，而是指着另一个对我们不可见的天，之后我们所见的这个天被称为\Quote{穹苍}，并且水的一半被提取上去，以供雨水、阵雨和露水给人类。另一半水留在地上，供河流、泉源和海洋之用。水曾覆盖全地，特别是低洼之处，然后天主藉着祂的圣言，使水聚在一处，使干地显露出来，这干地先前是看不见的。地成为可见后，仍是混沌的。于是天主用各种草木、种子和植物来塑造并装饰它。\par

% structure: p0060 source=h2 target=subsection reason=relative-heading-rank; base=h2

\subsection{第十四章 世界比作海}

再者，请思考它们的多样性、各异之美与繁多，以及如何借着它们显明复活，作为将来众人复活之预像。试问，有谁思及此，能不惊叹无花果树从无花果种子生出，或极其微小的种子竟能长成参天大树？我们说，世界犹如大海。就如大海若没有江河泉源涌入滋养，早已因咸性而枯竭；同样，世界若没有天主的法律和先知涌流而出甘饴、怜悯、正义和天主圣洁诫命的教训，也早已因其中充斥的邪恶和罪恶而灭亡。又如海中有岛屿，有些可居住、水源充足、果实丰饶，且有港湾和港口供受风暴者避难——同样，天主为那受罪恶驱迫与颠簸的世界，赐予了集会——我们是指圣教会——其中存有真理的教训，如同港湾中良好的锚地；凡渴望得救、热爱真理、但愿逃避天主的义怒与审判者，都投奔其中。又如另有岛屿，岩石累累、无水、荒芜、野兽出没、不可居住，只危害航海者和受风暴者，船只在那里遇难，漂流其上者丧亡——同样也有谬误的教训——我是指异端——它们毁灭接近它们的人。因为他们不受真理之言的引导；而如同海盗，装满船只后，将其驶至前述之地，以便掠夺：同样，凡偏离真理的人，也因自己的谬误而完全毁灭。\par

% structure: p0062 source=h2 target=subsection reason=relative-heading-rank; base=h2

\subsection{第十五章 论第四日}

第四日，光体被造成；因为天主预知虚妄哲学家的愚妄，他们要说地上生长的万物是从天体产生的，从而排除天主。因此，为使真理显而易见，植物和种子先于天体被造，因为后生的不能产生先有的。这些含有伟大奥秘的预像与样式。因为太阳是天主的预像，月亮是人的预像。就如太阳在能力和荣耀上远超月亮，同样天主远超于人。又如太阳永远圆满，永不减少，同样天主始终完美，充满一切能力、理解、智慧、不朽及一切美善。但月亮每月亏缺，仿佛死亡，为人预像；然后它再生，渐圆，作为将来复活的预像。同样，光体前三日是天主圣三——天主、祂的圣言和祂的智慧——的预像。第四日则是需要光的人之预像，如此便有天主、圣言、智慧、人。因此，光体在第四日被造成。星宿的排列也含有义人和虔诚者、以及遵守天主法律和诫命者的秩序与安排的预像。因为璀璨明亮的星宿是先知们的仿效，因此它们固定不移，不倾斜，不转换位置。那些次等明亮的，是义人百姓的预像。而那些变换位置、从一处逃往另一处的，即被称为行星的，也是那偏离天主、离弃祂的法律和诫命之人的预像。\par

% structure: p0064 source=h2 target=subsection reason=relative-heading-rank; base=h2

\subsection{第十六章 论第五日}

第五日，出于水的活物被造，借此也显明天主在这些事物中的多样智慧；因为谁能数算其众多和非常多样的种类？再者，出于水的事物蒙天主祝福，这也为预兆：凡来到真理、藉水和重生之洗重生、并从天主领受祝福的人，注定要接受悔改和罪过的赦免。但深海巨兽和猛禽是贪婪之人和违法者的相似。因为鱼和飞禽本属同一本性——有些持守本性，不加害比自己弱小的，遵守天主的法律，以地上的种子为食；另一些却违反天主的法律，吃肉，加害比自己弱小的。同样，义人遵守天主的法律，不咬人、不害人，圣洁公义地生活。但强盗、凶手和无神之人如同深海巨兽、野兽和猛禽；因为他们实质上吞噬比自己弱小者。因此，鱼类和爬虫的种族，虽分享了天主的祝福，却未获得特别显著的特性。\par

% structure: p0066 source=h2 target=subsection reason=relative-heading-rank; base=h2

\subsection{第十七章 论第六日}

在第六天，天主造了四足动物、野兽和地上的爬虫，未向它们降福，却为人类保留祂的降福，祂要在第六天造人。四足动物和野兽也被造为一些人的预像，这些人既不认识也不敬拜天主，只思念属地的事，且不悔改。因为那些远离不义、活出正义的人，在精神上如鸟向上飞翔，思念天上的事，中悦天主的旨意。但那些不认识也不敬拜天主的人，就像有翅膀却不能飞、不能翱翔到天主高处的鸟。同样，这样的人虽被称为人，却被罪压制，思念卑下和属地的事。动物被称为野兽（\OriginalTerm{野兽}{\GreekText{θηρία}}），因为它们被捕获（\OriginalTerm{被捕获}{\GreekText{θηρεύεσθαι}}），并非它们起初被造为邪恶或有毒的——因为天主所造的没有一样是邪恶的，但一切都是好的，甚至非常好——而是人所涉及的罪给它们带来了邪恶。因为当人犯罪时，它们也与他一同犯罪。正如如果家主自己行为端正，仆人也必然品行良好；但如果家主犯罪，仆人也与他一同犯罪。同样地，在人犯罪的情况下，他作为主人，所有隶属他的都与他一同犯罪。因此，当人再度回归他本性的状态，不再作恶时，那些动物也将恢复它们原本的温顺。\par

% structure: p0068 source=h2 target=subsection reason=relative-heading-rank; base=h2

\subsection{第十八章 人的创造}

至于关于人的创造，人自己无法解释，尽管圣经对此有简洁的记述。因为当天主说：\Quote{让我们按我们的形象、按我们的样式造人}时，祂首先暗示了人的尊贵。因为天主藉祂的圣言造了万有，并将它们都视为次要作品，唯有人的创造被视为配得上祂亲手的工作。此外，天主似乎需要帮助，说道：\Quote{让我们按我们的形象、按我们的样式造人。}但祂不是对别人，而是对祂自己的圣言和智慧说：\Quote{让我们造。}当祂造了人并降福他，使他增长并充满大地时，祂将万有置于他的管辖之下，供他使用；祂从一开始就规定他应从地上的果实、种子、菜蔬和橡子中获取营养，同时规定动物具有与人类似的习性，使它们也可以吃地上的种子。\par

% structure: p0070 source=h2 target=subsection reason=relative-heading-rank; base=h2

\subsection{第十九章 人被安置在乐园里}

天主既然在第六天完成了天地、海洋及其中的万有，便在第七天安息了祂所做的一切工作。然后圣经用这些话做了总结：\Quote{这是天地创造的记录：在主造天地之日，田野的每一种青草在被造之前，田间每一种菜蔬在生长之前。因为天主还没有降雨在地上，也没有人耕种土地。}\BibleRef{创 2:4}借此祂启示我们：那时全地靠着一口神圣的泉水灌溉，不需要人耕种；大地因天主的命令自发地产生万物，以免人因耕种而疲劳。但为了使人的创造得以明朗，以免在人中间似乎存在一个无法解决的问题——因为天主曾说过\Quote{让我们造人}，而祂的被造还没有明确叙述——圣经教导我们说：\Quote{有一道泉水从地上涌出，滋润全地面；天主用地上的尘土造人，将生命的气息吹入他的面中，人就成了有灵的生命。}因此大多数人称灵魂为不朽的。在造人之后，天主在东方之地中为他拣选了一个区域，在光方面卓越，因非常明亮的大气而光辉灿烂，[丰富的]最美好的植物，并将人安置在其中。\par

% structure: p0072 source=h2 target=subsection reason=relative-heading-rank; base=h2

\subsection{第二十章 圣经关于乐园的记述}

圣经这样记述神圣的历史：\BibleQuote{天主在东方、在\OriginalTerm{伊甸}{Eden}栽植了一个\OriginalTerm{乐园}{Paradise}，就把祂所形成的人安置在那里。天主使地面生出各种好看好吃的树，在乐园中央有\OriginalTerm{生命树}{tree of life}和\OriginalTerm{知善恶树}{tree of the knowledge of good and evil}。有一条河从伊甸流出浇灌乐园，并由那里分为四支：第一支名叫丕雄，环流\Place{哈腓拉}全境，此地出产黄金，黄金很好，又产芬芳胶脂和玛瑙；第二支河名叫基红，环流\Place{雇士}全境；第三支河名叫底格里，流向\Place{叙利亚}；第四支河就是\Place{幼发拉的}。主天主将人安置在乐园中，叫他耕种看守。天主命令\Person{亚当}说：乐园中各树上的果子你都可以吃，只有知善恶树上的果子你不可吃，因为你吃的日子必定死。主天主又说：人单独不好，让我们为他造一个相称的助手。于是主天主用土造成田野诸兽和天空诸鸟，都领到亚当前看他如何称呼；凡亚当给每个生物起的名称，就成了它的名字。亚当给一切牲畜、空中的鸟和田野的诸兽都起了名；但是亚当没有找到相称的助手。于是天主使亚当沉睡，他睡着了，就取出一根肋骨，补上肉。主天主用那根肋骨造成一个女人，领到亚当前。亚当说：这才是我的骨中之骨，肉中之肉；她应称为女人，因为她是从男人取出来的。为此男人要离开父母，依恋他的妻子，二人成为一体。当时亚当和他的妻子都赤身露体，并不害羞。}\par

% structure: p0074 source=h2 target=subsection reason=relative-heading-rank; base=h2

\subsection{第二十一章 论人的堕落}

\BibleQuote{现在，\OriginalTerm{蛇}{serpent}是田野诸兽中最狡猾的。蛇对女人说：天主真说过，你们不可吃园中任何树上的果子么？女人对蛇说：我们可以吃园中树上的果子，只有园中央那棵树上的果子，天主说过，你们不可吃，也不可摸，怕你们会死。蛇对女人说：你们决不会死！因为天主知道，你们吃的那天，你们的眼就会开，你们就会像神一样，知道善恶。女人看见那棵树实在好吃，又悦目，且能使人增长智慧，就摘下果子吃了，又给了与她一起的丈夫，他也吃了。于是两人的眼就开了，发觉自己赤身露体，就把无花果叶缀在一起，编了裙子。他们听见主天主在园中趁晚凉散步的声音，\Person{亚当}和他的妻子就躲藏在园中的树木中，逃避主天主的面。主天主呼唤亚当说：你在哪里？他答说：我听见你在园中的声音，就害怕起来，因为我赤身露体，就藏了起来。祂说：谁告诉你赤身露体？莫非你吃了我禁止你吃的果子？亚当说：是你给我作伴的那个女人，把那树上的果子给了我，我就吃了。主天主对女人说：你做了什么事？女人说：是蛇诱骗了我，我就吃了。主天主对蛇说：你做了这事，在一切牲畜和田野的野兽中，你当受咒骂；用你的肚皮爬行，一生以尘土为食。我要把仇恨放在你和女人之间，放在你的后裔和她的后裔之间；她的后裔要踏碎你的头颅，你要伤害他的脚跟。对女人祂说：我要加重你怀孕的苦楚，在痛苦中生产儿女；你要依恋你的丈夫，他要管辖你。对亚当祂说：因为你听从了你妻子的话，吃了我禁止你吃的果子，为了你的缘故，地成了可咒骂的；你一生日日劳苦才能得到吃食。地要给你生出荆棘和蒺藜，你要吃田间的野菜；你必须汗流满面，才有饭吃，直到你归于土，因为你来自土；你原是尘土，还要归于尘土。}这就是圣经关于人类历史和乐园的记述。\par

% structure: p0076 source=h2 target=subsection reason=relative-heading-rank; base=h2

\subsection{第二十二章 为何说天主行走}

你会对我说：\Quote{你说过天主不应被局限在一个地方，现在你怎么又说祂在乐园中行走呢？}请听我说。万有的天主和圣父确实不能被局限，也不在任何地方，因为没有祂安息之所；但祂的圣言——祂藉之创造万物的那位，即祂的能力和智慧——取了万有的父和主的位格，以天主的位格来到园中，与亚当交谈。因为圣经本身教导我们，亚当说他听到了声音。但这声音若非天主的圣言（祂也是天主的圣子），又是什么呢？并非如诗人和神话作者所谈论的那种从交合中产生的众神之子，而是如真理所阐明的：圣言永远存在，居住在天主的心内。在万物存在之前，天主就拥有祂作为谋士，即祂自己的理智和思想。但当天主愿意成就祂所决定的一切时，祂就生了这圣言，发出之声，一切受造物的首生者，并非祂自己失去了圣言（理性），而是生了理性，并常与祂的理性交谈。因此，圣经和所有受圣神感动的人都教导我们，其中一位若望说：\Quote{在起初已有圣言，圣言与天主同在，}\BibleRef{若 1:1}表明起初天主独自存在，圣言在祂内。然后他说：\Quote{圣言就是天主；万物是藉着祂而造成的；凡受造的，没有一样不是由祂造成的。}圣言既是天主，从天主自然地生出，每当宇宙之父愿意时，祂就派遣圣言到任何地方；圣言来临时，既被听见又被看见，由祂所派遣，并出现在一个地方。\par

% structure: p0078 source=h2 target=subsection reason=relative-heading-rank; base=h2

\subsection{第23章　创世记记述的真实性}

因此，天主在第六天造了人，并在第七天后启示了这创造，那时祂也造了乐园，使人能在一个更美好且明显更优越的地方。而这是真实的，事实本身证明了这一点。因为谁能看不见，妇女在分娩时所受的痛苦和随后她们对劳苦的遗忘，都是为成就天主的话，使人类繁衍增多呢？我们不是也看见了对蛇的判决——它可憎地以腹部爬行并以尘土为食——使我们也有这个作为先前所说之事的证据吗？\par

% structure: p0080 source=h2 target=subsection reason=relative-heading-rank; base=h2

\subsection{第24章　乐园的美丽}

于是，天主使地上长出各种美观悦目或可作食物的树。因为起初只有第三天所产生的东西——植物、种子和菜蔬；但乐园中的东西具有更优越的可爱和美丽，因为其中的植物据说是天主所栽种的。至于其余植物，世界确实有类似的植物；但那两棵树——生命树和知善恶树——其余的地上都没有，只有乐园才有。而乐园是地，栽种在地上，圣经说：\BibleRef{创 2:8}\Quote{上主天主在伊甸东部种植了一个乐园，就将他所造的人安置在里面。上主天主使地面生出各种好看好吃的树。}因此，藉着\Quote{地面}和\Quote{东部}这些词语，圣经清楚地教导我们，乐园是在这个天之下，东和地也在其下。希伯来语\OriginalTerm{厄登}{Eden}意为\Quote{快乐}。并且预示有一条河从厄登流出灌溉乐园，然后分为四支；其中两条称为\OriginalTerm{丕雄}{Pison}和\OriginalTerm{基红}{Gihon}灌溉东部地区，尤其是基红环绕整个雇士地，据说它在埃及以尼罗河之名重现。另外两条河我们显然可以辨认——称为底格里斯河和\OriginalTerm{幼发拉的河}{Euphrates}——因为它们邻近我们的地区。天主将人安置在乐园中，如前所述，为耕种和看守，命令他吃所有树上的果实——显然也包括生命树；但唯独知善恶树上的果实祂命令他不可尝。天主将他从所出的地上迁到乐园，赐给他进步的方法，使他成熟、完善，甚至被称为神，这样他就能拥有不朽而升天。因为人被造成一种中间本性，既非完全可朽，也非全然不朽，而是可能两者；同样，乐园这个地方在美丽方面被造为介于天地之间。藉着\Quote{耕种}一词，所暗示的劳动无非是遵守天主的命令，以免他因违命而自取灭亡，正如他确实因罪自灭了。\par

% structure: p0082 source=h2 target=subsection reason=relative-heading-rank; base=h2

\subsection{第25章　天主禁止人吃知善恶树上的果实是公义的}

知识树本身是好的，它的果子也是好的。因为不是树，像有些人所想的那样，而是不服从，其中有死亡。因为果子里没有别的，只有知识；但当人明智地使用知识时，知识是好的。然而，亚当在年龄上尚属婴孩，因此还不能相称地领受知识。因为如今，当一个孩子出生时，他也不能立刻吃面包，而是先用奶喂养，然后随着年岁的增长，才进步到固体食物。对亚当也是如此；因为天主并非像有些人所猜测的那样，出于吝啬而不许他吃知识树。相反，天主也要考验他，看他是否服从祂的诫命。同时，天主愿意人，尚为婴孩，在一段时期内保持纯朴和真诚。因为无论在神面前还是人面前，以纯朴和真诚顺服父母，这是神圣的。但如果儿女应当顺服父母，何况万有的天主和圣父呢？此外，婴孩在年幼时智慧超过年龄是不合宜的；因为正如人在身量上有顺序地增长，在智慧上也是如此。但是，正如当一条法律禁止某事，而有人不遵守时，显然不是法律导致惩罚，而是不服从和违犯——因为父亲有时命令自己的孩子戒绝某些事，当孩子不服从父亲的命令时，他因不服从而受鞭打和惩罚；在这种情况下，行为本身并不是鞭打的原因，而是不服从为不服从者招致惩罚——同样，对于第一个人，不服从导致他被逐出乐园。因此，并非知识树本身有任何邪恶，而是由于不服从，人如同从泉源中汲引出劳苦、痛苦、忧伤，并最终被死亡所攫获。\par

% structure: p0084 source=h2 target=subsection reason=relative-heading-rank; base=h2

\subsection{第二十六章：天主将人逐出乐园的恩慈。}

天主在此事上向人显示了极大的恩慈：祂不允许人永远停留在罪恶中；而是，可以说，借着一种放逐，把他逐出乐园，使他能在指定时间内借着惩罚赎罪，受到管教，然后得以恢复。因此，当人在这个世界被造时，在\BibleBook{}{创世纪}中奥秘地记载，好像他曾两次被安置在乐园中；这样，第一次在他被安置时就应验了，第二次将在复活和审判之后应验。因为正如一个器皿，在制作时发现有瑕疵，就被重新塑造或重做，使之成为新的、完整的；同样，人借着死亡也是如此。因为人某种程度上被拆毁，为的是在复活中完整地兴起；我是指纯洁、正义且不朽坏的。至于天主的呼唤和说：\BibleQuote{亚当，你在哪里？}天主这样做，并非不知道此事；而是由于祂的容忍，给亚当悔改和告解的机会。\par

% structure: p0086 source=h2 target=subsection reason=relative-heading-rank; base=h2

\subsection{第二十七章：人的本性。}

但有人会对我们说：人是本性可朽的吗？当然不是。那么，他是不朽的吗？我们也不肯定这一点。但有人会说：那么他什么也不是吗？这也没有击中要害。他本性既非可朽，也非不朽。因为如果天主从一开始就使他成为不朽的，就会使他成为天主。再者，如果祂使他成为可朽的，天主就似乎成为他死亡的原因。因此，天主既没有使他成为不朽的，也没有使他成为可朽的，而是如我们上面所说，使他能够承受两者；这样，如果他倾向于不朽之事，遵守天主的诫命，他就从祂那里得到不朽作为赏报，并成为天主；但另一方面，如果他转向死亡之事，不服从天主，他就成为自己死亡的原因。因为天主造了自由的人，并赋予他自主权。因此，人因疏忽和不服从而给自己带来的，现在当天主因自己的仁慈和怜悯而赐予服从祂的人作为恩赐。因为正如人因不服从而将死亡招致自身；同样，服从天主的旨意，凡愿意的人都能为自己获得永生。因为天主赐给了我们法律和神圣的诫命；凡遵守这些诫命的人都能得救，获得复活，并继承不朽。\par

% structure: p0088 source=h2 target=subsection reason=relative-heading-rank; base=h2

\subsection{第二十八章：夏娃为何由亚当的肋骨形成。}

\Person{亚当}被逐出乐园后，在这样的状态下认识了自己的妻子\Person{厄娃}，她是\Person{天主}从\Person{亚当}的肋骨中为他形成的一个妻子。天主这样做，并非因为他不能单独造出妻子，而是天主预知人会呼求许多神。天主有这预知，并知道通过蛇，错误将引入许多并不存在的神——因为只有一位神，即使那时错误也正试图散布众多神，说：\Quote{你们将如同神一样；}——因此，免得有人以为一位神造了男人，另一位神造了女人，所以天主造了他们两者；而且天主将女人与男人一同造成，不仅是为了显明天主独一治理的奥迹，也是为了使他们之间的相互爱慕更加深厚。因此\Person{亚当}对\Person{厄娃}说：\Quote{这才是我的骨中之骨，肉中之肉。}此外，他还预言说：\Quote{为此，人要离开自己的父亲和母亲，依附自己的妻子，两人成为一体；}这个预言在我们自己身上也得应验。因为谁合法结婚，不轻视母亲和父亲，以及整个亲属关系和家庭，而依附并与自己的妻子成为一体，深情地偏爱她呢？因此，为了妻子，有些人甚至常常赴死。\Person{厄娃}，因她起初被蛇欺骗而成为罪的肇始者，那邪恶的魔鬼，也称为\Person{撒殚}，它当时通过蛇对她说话，并且至今仍在那被它附身的人身上工作，它却以厄娃之名呼求。它被称为\Quote{魔鬼}和\Quote{龙}，因为它\OriginalTerm{背弃}{\GreekText{ἀποδεδρακέναι}}了天主。起初它是一位天使。关于它的历史有许多可说的；因此我目前省略其叙述，因为我已在另一处记述了它。\par

% structure: p0090 source=h2 target=subsection reason=relative-heading-rank; base=h2

\subsection{第二十九章 加音的罪行}

当\Person{亚当}认识了他的妻子\Person{厄娃}，她便怀孕生了一个儿子，取名\Person{加音}；她说：\Quote{我赖天主获得了一个人。}之后她又生了第二个儿子，取名\Person{亚伯尔}，\Quote{亚伯尔成了牧羊者，加音则耕田。} \BibleRef{创 4:1} 他们的历史有非常详尽的叙述，甚至是详细的解释；因此那部名为\WorkTitle{创世纪}的书本身能更准确地告知那些渴望了解他们故事的人。当\Person{撒殚}看见\Person{亚当}和他的妻子不仅仍活着，而且还生儿育女——因未能置他们于死地而心怀恶意——当他看见\Person{亚伯尔}为天主所喜悦时，他便煽动那叫\Person{加音}的兄弟的心，使他杀害了自己的兄弟\Person{亚伯尔}。如此，死亡便在这个世界有了开端，并进入每一人类种族，直到今日。但\Person{天主}是慈悲的，愿意给\Person{加音}，如同给\Person{亚当}一样，一个悔改和认罪的机会，便说：\Quote{你的弟弟\Person{亚伯尔}在哪里？}但\Person{加音}悖逆地回答天主说：\Quote{我不知道；难道我是我弟弟的看守者吗？}天主因此向他发怒说：\Quote{你做了什么事？你弟弟的血由地上向我喊冤。地开口接受了你兄弟的血从你手中流下。你在地上要呻吟颤抖。}\par

% structure: p0092 source=h2 target=subsection reason=relative-heading-rank; base=h2

\subsection{第三十章 加音的家族及其发明}

\Person{加音}自己也有一个儿子，名叫\Person{哈诺客}；他建造了一座城，以他儿子的名字命名为\Place{哈诺客}。从那时起，开始了城市的建造，这发生在洪水之前；不像\Person{荷马}错误地所说：\par

\Quote{那时人们尚未建造城市。}\par

\Person{哈诺客}生了一个儿子，名叫\Person{依辣得}；\Person{依辣得}生\Person{默胡雅耳}；\Person{默胡雅耳}生\Person{默突沙耳}；\Person{默突沙耳}生\Person{拉默客}。\Person{拉默客}娶了两个妻子，名叫\Person{阿达}和\Person{漆拉}。那时开始了多妻制，也开始了音乐。因为\Person{拉默客}有三个儿子：\Person{雅巴耳}、\Person{犹巴耳}、\Person{突巴耳}。\Person{雅巴耳}成了畜牧者，住在帐棚中；\Person{犹巴耳}是传授琴瑟的祖师；\Person{突巴耳}成了铜匠铁匠。以上是\Person{加音}的后裔；其余的他这一支的后裔因他杀害兄弟而被遗忘。代替\Person{亚伯尔}，\Person{天主}赐予\Person{厄娃}怀孕生子，取名\Person{舍特}；人类其余的部分从\Person{舍特}延续至今。对于渴望知晓所有世代的人，借助圣经很容易解释。因为，正如我们已经提到过的，这个关于人类谱系顺序的主题，在我们另一篇讲论中，在\WorkTitle{历史书}的第一卷，已有部分论述。这些事都是圣神教导我们的，祂通过\Person{梅瑟}和其他先知说话，因此我们虔敬之人的著作比所有作家和诗人都更古老，并且被证明更为真实。此外，关于音乐，有些人虚构说\Person{阿波罗}是发明者，另一些人说\Person{俄耳甫斯}从鸟的悦耳歌声中发现了音乐艺术。他们的故事被证明是空洞和虚妄的，因为这些发明者生活在洪水之后许多年。至于\Person{诺厄}，有些人称他为\Person{丢卡利翁}，我们在前述的书中已作解释，如果你愿意，你可以阅读那卷书。\par

% structure: p0096 source=h2 target=subsection reason=relative-heading-rank; base=h2

\subsection{第三十一章 洪水后的历史}

洪水过后，城市和诸王又有了新的开端，情形如下：——第一座城是\Place{巴比伦}、\Place{厄勒客}、\Place{阿加得}和\Place{加耳乃}，都在\Place{史纳尔}地。他们的王名叫\Person{尼默洛得}。从这些人中出来的有\Person{亚述}，亚述人也由此得名。\Person{尼默洛得}建造了\Place{尼尼微}、\Place{勒曷波特}、\Place{加拉赫}和\Place{勒森}，在\Place{尼尼微}与\Place{加拉赫}之间；\Place{尼尼微}成了一座极大的城。\Person{诺厄}的儿子\Person{闪}的另一个儿子，名叫\Person{米兹辣因}，生了\Person{路丁人}、\Person{阿纳明人}、\Person{肋哈宾人}、\Person{纳夫突欣人}、\Person{帕特洛斯人}和\Person{加斯路欣人}，从他们中出来了\Person{培肋舍特人}。关于\Person{诺厄}的三个儿子，他们的死亡和谱系，我们在上述书中已给出了简明的记录。但现在我们要提及其余的事实，关于城市和诸王，以及当人们有同一语言和同一口音时所发生的事。在语言分裂之前，上述的城市就已存在。但当人们将要分散时，他们凭自己的判断商议，而不是出于天主的推动，要建造一座城和一座塔，塔顶可通天上，好为自己扬名。因此，既然他们违抗天主的旨意，胆敢从事一项大工程，天主就摧毁了他们的城，推倒了他们的塔。从那时起，他混乱了人的语言，赐给各人不同的方言。同样，\Person{西比拉}也曾发言，说她得知愤怒将临于世界。她说：——\par

\Quote{当伟大天主的威胁得以实现，\newline{}他以这些威胁警告世人，当初在\newline{}亚述之地他们建造了一座塔，\newline{}众人都说同一种语言，并想上升\newline{}直到他们攀上星辰的天空，\newline{}那时永生者兴起了一阵大风，\newline{}并将强大的必然性加于他们；\newline{}因为当那风推倒了那座大塔，\newline{}在人群中就兴起了激烈的纷争和仇恨。\newline{}一种语言变成了多种方言，\newline{}大地充满了各种部落和诸王。}\par

如此等等。这些事发生在\Place{加色丁}人的地。在\Place{客纳罕}地有一座城，名叫\Place{哈兰}。在那些日子，\Person{法郎}（埃及人也叫他\Person{乃苛}）是埃及的第一个王，之后诸王相继。在\Place{史纳尔}地，在那些称为\Place{加色丁}人当中，第一个王是\Person{阿黎约客}，接着他的是\Person{厄拉撒尔}，之后是\Place{厄蓝}王\Person{革多尔老默尔}，再后是称为亚述的各民族之王\Person{提达耳}。还有五座城在\Person{诺厄}的儿子\Place{含}的境内：第一座叫\Place{索多玛}，然后是\Place{哈摩辣}、\Place{阿德玛}、\Place{则波殷}和\Place{贝拉}（也称为\Place{左哈尔}）。他们的王的名字如下：\Place{索多玛}王\Person{贝辣}，\Place{哈摩辣}王\Person{彼尔沙}，\Place{阿德玛}王\Person{史纳布}，\Place{则波殷}王\Person{舍默贝尔}，\Place{贝拉}（即\Place{克法拉克}）王\Person{贝拉}。这些王服事亚述王\Person{革多尔老默尔}十二年，到了十三年他们背叛了\Person{革多尔老默尔}；于是在那时发生了四个亚述王与五个王的战争。这是大地上战争的首次开端；他们摧毁了\Place{卡尔纳殷}的巨人，以及与他们在城中的那些强大的民族，和称为\Place{色依尔}的山地上的\Person{曷黎人}，直到靠近旷野的\Place{帕兰}平原。在那时，在\Place{撒冷}（即现在的\Place{耶路撒冷}）城中有一位正义的王，名叫\Person{默基瑟德}。他是至高者天主所有司铎中的第一位司铎，从此上述的城\Place{耶路撒冷}得名。从那时起，司铎在全地出现。在他之后，\Person{阿彼默肋客}在\Place{革辣尔}统治；之后是另一个\Person{阿彼默肋客}。然后\Person{厄斐龙}（被称为赫特人）统治。这些是古时诸王的名字。其余的亚述诸王，在许多年的间隔中，均被静默忽略而不录，所有作者都叙述我们近日的事。这些是亚述的王：\Person{提革拉特丕肋色尔}，之后是\Person{沙耳玛乃色尔}，然后是\Person{散乃黑黎布}；以及\Person{阿德辣默肋客}（雇士人），他也统治过埃及，是他的三头统治之一；——不过这些事与我们的书相比，则是相当晚近的。\par

% structure: p0100 source=h2 target=subsection reason=relative-heading-rank; base=h2

\subsection{第三十二章 人类如何分散}

因此，爱好学问和古事的人可以由此理解历史，并看到那些我们离开圣先知们所叙述的事都是晚近的。因为起初在\Place{阿拉伯}和\Place{加色丁}地只有少数人，然而，在他们的语言分裂之后，他们逐渐开始增多并散布到全地；有的趋向东方去居住，有的趋向大洲的部分，有的趋向北方，以至于延伸到北极地区的\Place{不列颠}。还有的到了\Place{客纳罕}地（即\Place{犹太}）和\Place{腓尼基}，以及\Place{厄提约丕雅}地区、\Place{埃及}、\Place{利比亚}，和称为\Quote{灼热}的地方，以及伸展向西的部分；其余的人则去了沿海的地方、\Place{潘非利亚}、\Place{亚细亚}、\Place{希腊}、\Place{马其顿}，此外还有\Place{意大利}、整个称为\Place{高卢}的地方、\Place{西班牙}和\Place{日耳曼}；所以现在全地就这样充满了居民。因为人类占据世界起初分为三个部分——东方、南方和西方；后来当人变得非常多时，大地剩余的部分也住了人。但作者们，不知道这些事，竟贸然主张世界是球形的，并将之比作一个立方体。但他们对这些事怎能说真话呢？他们连世界的创造和它的人口都不知道。人逐渐增多并在地上繁殖，如我们已说过的，海中的岛屿和其余的国家也住了人。\par

% structure: p0102 source=h2 target=subsection reason=relative-heading-rank; base=h2

\subsection{第三十三章 世俗历史未记载这些事}

那么，那些被称为智者、诗人和史家的人，他们自身出生晚得多，又引进许多在城邦之后多年才出生、比君王、民族和战争更近代的神祇，有谁能真实地告诉我们这些事呢？因为他们本应提及所有事件，甚至那些洪水之前发生的事，包括世界的创造、人的形成以及全部事件的延续。埃及或迦勒底的先知，以及其他作者，如果至少他们藉着神圣纯洁的神灵发言，并在他们所说的一切中讲真话，就应该能够准确地讲述；而且他们不仅应宣布过去或现在的事，也应宣布将来临到世界的事。因此，这便证明所有其他人都错了；唯独我们基督徒拥有真理，因为我们是由圣神教导的，祂曾藉圣先知发言，并预言了一切。\par

% structure: p0104 source=h2 target=subsection reason=relative-heading-rank; base=h2

\subsection{第三十四章 先知们颁令圣洁生活。}

此外，但愿你以善意的心态探究神圣的事——我指的是先知们所说的事——以便通过将我们所说的与别人的言论相比较，你能够发现真理。我们已经从他们自己编纂的历史中表明，那些被称为神祇的名字，正是他们中间曾经生活过的人的名字，正如我们上面所证明的。直到今日，他们的像每天被塑造，即偶像，\Quote{人手之工}。愚昧的大众事奉这些偶像，却拒绝万物的创造者、塑造者以及一切生命气息的滋养者，因他们从祖先那里接受的无知传统的欺骗而相信虚妄的教理。然而，天主，宇宙的圣父和创造者，并没有放弃人类，而是赐下法律，派遣圣先知来宣告和教导人类，使我们每个人都能醒悟并明白惟有一天主。他们也教导我们远离非法的偶像崇拜、奸淫、杀人、邪淫、偷盗、贪婪、虚誓、愤怒以及一切不节制和不洁；并且，一个人不愿他人对自己做的事，也不应加诸他人；这样，行义的人将逃脱永罚，并被认为配得天主所赐的永生。\par

% structure: p0106 source=h2 target=subsection reason=relative-heading-rank; base=h2

\subsection{第三十五章 先知书的训诫。}

因此，神圣的法律不仅禁止敬拜偶像，也禁止敬拜天上的天体、太阳、月亮或其他星辰；是的，也不可敬拜天、地、海、泉源或河流，而必须以圣洁的心和真诚的意向事奉那永生且真实的天主，祂也是宇宙的创造者。为此神圣的法律说：\BibleQuote{不可奸淫；不可偷盗；不可作假见证；不可贪恋你邻人的妻子。}先知们也如此说。撒罗满确实教导我们连眼目的一转也不可犯罪，说：\BibleQuote{愿你眼睛向前直视，愿你眼睑直看前方。} \BibleRef{箴 4:25} 梅瑟本人也是一位先知，论及天主的独一统治说：\BibleQuote{你的天主是那设立天、形成地、以手带出天上万军的那位；祂并没有把这些摆在你们面前，叫你们去随从它们。} \BibleRef{申 4:19} 依撒意亚本人也说：\BibleQuote{上主天主如此说：祂设立天、奠定地、赐气息给地上的人民，并赐灵给行在其间的人。这是你的上主天主。} \BibleRef{依 42:5} 又藉着祂说：\BibleQuote{我造了地，也造了地上的人；我亲手设立了天。} \BibleRef{依 45:12} 在另一章中：\BibleQuote{这是你的天主，祂创造了地极；祂不饥不倦，祂的智慧无法测度。} \BibleRef{依 40:28} 同样，耶肋米亚说：\BibleQuote{祂以能力造了地，以智慧建立了世界，以明达展开了天，并有一团水在天上；祂使云彩从地极上升，祂造闪电带雨，从祂的宝库中带出风来。} \BibleRef{耶 10:12} 由此可见，所有先知如何一致而和谐地发言，藉着同一个圣神发言，论及天主的统一、世界的创造和人的形成。此外，他们极其痛苦，哀叹无神的人类，并责备那些自以为有智慧的人的错误和心硬。耶肋米亚确实说过：\BibleQuote{每个人都因不认识祂而愚拙地偏离；每个工匠都因他的雕刻偶像而羞愧；银匠徒然铸造可铸的像；它们里面没有气息；在他们受罚的日子，他们必灭亡。} \BibleRef{耶 51:17} 达味同样说：\BibleQuote{他们都败坏了，行了可憎的事；没有一个行善的，连一个也没有；他们都离弃了正道，一同变成无益的。} 哈巴谷也说：\BibleQuote{雕刻的偶像有何益处呢？雕刻者雕刻了一个虚谎的像。祸哉，那对石头说：醒来；对木头说：起来。} \BibleRef{哈 2:18} 其他先知也同样说了真理。我何必列举众多的先知呢？他们人数众多，说了成千上万彼此一致而和谐的话。因为凡愿意的人，可以藉着阅读他们所说的，准确地理解真理，不再被意见和无益的辛劳所迷惑。这些我们已提及的先知，就是希伯来人中的先知——不识字、是牧人、未受教育。\par

% structure: p0108 source=h2 target=subsection reason=relative-heading-rank; base=h2

\subsection{第三十六章 西比尔的预言。}

西比尔，在希腊人和其他民族中是一位女先知，在她预言的开端，责备人类说：——\par

\Quote{你为何如此轻易就被高举，\newline{}又为何如此不顾生命的终结，\newline{}你们这些属血肉的凡人，不过是虚无？\newline{}你们岂不战栗，也不敬畏至高的天主？\newline{}你们的监护者，那知悉、洞悉万有的，\newline{}祂永远保守祂亲手所造的第一批，\newline{}将祂的甘饴圣神注入祂的一切化工，\newline{}并将祂赐予世人作为向导。\newline{}只有一位非受造的天主，\newline{}祂独掌王权，全能、至大，\newline{}无物能瞒过祂。祂看见一切，\newline{}自己却不被任何肉眼看见。\newline{}凡人岂能看见不死的天主，\newline{}或是血肉之眼，它们避开正午的光芒，\newline{}岂能仰望那居于诸天之上的祂？\newline{}那么，敬拜祂，那自有的天主，\newline{}那世界非受生的统治者，\newline{}祂自亘古就已存在，\newline{}并且将永远存续。\newline{}你们将为恶谋收割果实，\newline{}因为你们没有尊崇真神，\newline{}也没有向祂献上圣洁的百牛祭。\newline{}你们向居住在阴府中的献礼，\newline{}并向魔鬼献祭。\newline{}你们在癫狂和骄傲中行走；\newline{}离开正路，你们四处游荡，\newline{}失落在坑中和荆棘里。\newline{}愚昧的人啊，你们为何这样游荡？\newline{}停止你们在黑暗黑夜中的徒然徘徊；\newline{}为何追逐黑暗和永久的幽暗\newline{}当有福的光明为你们照耀？\newline{}看，祂是光明的——在祂内毫无污点。\newline{}那么，转离黑暗，仰望白日；\newline{}要有智慧，将智慧珍藏在你们心里。\newline{}有一位天主，祂兴起风雨、\newline{}地震、闪电、瘟疫、\newline{}饥荒、雪暴、冰雹，\newline{}以及降临我们悲苦一族的一切灾祸。\newline{}不仅这些，祂还赐予一切，\newline{}在天地间全能统治，\newline{}并且自亘古即自有。}\par

至于那些据说是被生出来的[神]，她说道：—\par

\Quote{若一切生者必死，\newline{}天主不能由凡人产生。\newline{}但唯有独一者，至上者，\newline{}祂创造了诸天及其众星，\newline{}太阳和月亮；同样丰饶的大地，\newline{}连同海洋的一切波涛，山丘，\newline{}泉源，和永远流淌的溪流；\newline{}祂也创造了无数\newline{}海洋生物，并永保存活\newline{}一切爬虫，地上的和海中的；\newline{}祂也造了飞鸟的歌唱之队，\newline{}它们以翅膀划破空气，以尖细的笛声\newline{}在清晨颤音鸣唱它们柔和清脆的歌声。\newline{}在山丘的深谷中祂安置了\newline{}野性的兽类；并向人\newline{}使一切牲畜驯服，让人成为\newline{}神形的肖像，统治万有，\newline{}并将许多不可理解之物\newline{}臣服于其权下。\newline{}因为凡人谁能知道这一切？\newline{}唯有在起初创造它们的祂知道，\newline{}祂是创造者，不朽坏的，\newline{}永远居住在至高之处；\newline{}祂向善者赐予极丰厚的赏报，\newline{}并向邪恶和不义的人\newline{}激起报复、愤怒、血战，\newline{}瘟疫，和许多流泪的悲伤。\newline{}徒然被高举的人啊——说为何\newline{}你如此彻底地毁了自己？\newline{}你敬拜兽类为神，难道不羞耻吗？\newline{}相信众神会偷走你的牲畜，\newline{}或需要你的器皿——这岂不是\newline{}癫狂最无益且愚昧的想法？\newline{}你们本该住在金色诸天，\newline{}却看见你们的神成为虫子的猎物，\newline{}和许多可憎不洁的受造物。\newline{}愚人啊！你们敬拜蛇、狗、猫，\newline{}鸟，和地上海中的爬虫，\newline{}人手所造之像、石雕，\newline{}以及路旁放置的垃圾堆。\newline{}所有这些，以及更多虚妄之物，你们事奉，\newline{}敬拜那些连提起都羞耻之物：\newline{}这些是引导虚妄之人走入歧途的神，\newline{}从他们口中流出致命毒液之流。\newline{}但是要向着那独有生命的祂，\newline{}生命、不死、永恒之光；\newline{}祂将喜悦注入人的生命之杯，\newline{}比最甜的蜜更甘甜——\newline{}向祂低头，唯独向祂，\newline{}行走在永恒平安的路上。\newline{}离弃祂，你们都转离正路，\newline{}并在疯狂的愚妄中，饮尽了\newline{}公义之杯，未掺杂、纯正、主宰、强烈；\newline{}你们不会再成为清醒的人，\newline{}不会恢复清醒的头脑，\newline{}认识你们的天主和君王，祂鉴察一切：\newline{}因此，烈火将降临你们，\newline{}你们将永远在火焰中日日焚烧，\newline{}永远为你们无用的神感到羞耻。\newline{}但那些敬拜永恒天主的人，\newline{}他们将继承永生，\newline{}居住在幸福盛开的领域，\newline{}享用来自星天之海的美物。}\par

这些事是真实、有益、正义且对所有人有益的，这是显然的。连诗人们也曾论及恶人的惩罚。\par

% structure: p0114 source=h2 target=subsection reason=relative-heading-rank; base=h2

\subsection{第三十七章 诗人们的见证}

恶人必按其所行受罚，这早已被一些诗人如同神谕般宣告，作为对他们自己和对恶人的见证，声明他们将受惩罚。\Person{埃斯库罗斯}说：—\par

\Quote{谁做了，也必须受。}\par

\Person{品达}自己说：—\par

\Quote{受苦紧随行为是合宜的。}\par

同样，\Person{欧里庇得斯}：—\par

\Quote{行为曾使你喜悦——受苦须忍受；\newline{}被擒的敌人必定痛苦。}\par

再次说：—\par

\Quote{敌人的痛苦是英雄的奖赏。}\par

同样，\Person{阿尔基洛科斯}：—\par

\Quote{有一事我知道，我永远认为它真实，\newline{}作恶者必遭受邪恶。}\par

至于天主洞悉万物，无物逃过祂的注意，但祂因长久忍耐而克制，直到祂要审判的时候——关于这点，\Person{狄奥尼修斯}也曾说：—\par

\Quote{正义之眼洞悉一切，\newline{}却似乎视而不见。}\par

而且天主的审判将要来临，灾祸将突然临到恶人——这同样由\Person{埃斯库罗斯}宣告说：—\par

命运之脚步迅疾，\newline{}无人能违反正义，\newline{}迟早会感到其严厉之手。\newline{}它与你同在，虽未闻、未见；\newline{}你拉下夜的帷幕于其间，\newline{}但连睡眠也不能提供屏障。\newline{}无论你眠或醒，它都与你同在；\newline{}若你外出上路，\newline{}你无法打破它静默严厉的守望。\newline{}它将跟随你，或横阻你的路；\newline{}甚至黑夜也无能力\newline{}将你隐藏以免复仇者的忿怒。\newline{}\Quote{为显示恶人，黑暗逃遁；\newline{}那么，若罪恶提供快乐或安逸，\newline{}哦，停住，想想有人看见！}\par

我们难道不能也引用\Person{西蒙尼德斯}吗？—\par

\Quote{对人来说，没有灾祸是无征兆降临的；\newline{}但天主以突如其来的手转变万物。}\par

\Person{欧里庇得斯}再次说：—\par

\Quote{恶人发达如沙基，\newline{}其族不久必倾颓；\newline{}时日宣告恶人罪。\newline{}}\par

\Person{Euripides}又说：—\par

\Quote{神明断非无判断，\newline{}但见人妄立誓言，\newline{}迫人强允定不宽。\newline{}}\par

\Person{Sophocles}又说：—\par

\Quote{行恶者必受恶报。}\par

天主必追究妄誓及一切恶事，他们自己几乎已预言此事。关于世界的焚毁，他们虽远较先知为晚，却或愿或不愿地依照先知说了话，是从法律和先知那里偷来的。诗人们证实了先知的见证。\par

% structure: p0138 source=h2 target=subsection reason=relative-heading-rank; base=h2

\subsection{第三十八章　希腊诗人和哲学家的教导与希伯来先知的教导相符}

但他们在先知之前或之后有什么关系呢？无论如何，他们确实说了与先知相符的话。关于世界的焚毁，先知\Person{玛拉基亚}预言说：\BibleQuote{看，主的日子来到，如燃烧的火炉，要消灭所有恶人。}\BibleRef{拉 4:1}\Person{依撒意亚}说：\BibleQuote{因为天主的忿怒如同猛烈的冰雹，如同奔腾的山洪。}\BibleRef{依 30:30}那么，女先知\Person{西比尔}和其他先知，乃至诗人和哲学家，都已清楚教导了关于正义、审判和惩罚，以及关于天主眷顾，即天主看顾我们，不仅看顾我们中的生者，也看顾死者；虽然他们是不情愿地说的，因为真理使他们信服。在先知的书中，\Person{撒罗满}论及死者说：\BibleQuote{你的肌肉必得医治，你的骨骸必受照顾。}\BibleRef{箴 3:8}同样，\Person{达味}也说：\Quote{你所折断的骨骸必要欢欣。}与这些说法一致的是\Person{提摩克勒斯}的话：—\par

\Quote{仁慈的天主怜悯死者。}\par

那些主张多神的作家最后也归到一神的道理，主张偶然的人也讲论天主眷顾，主张不受惩罚的人也承认有审判，否认死后有感觉的人也承认有感觉。因此，\Person{荷马}虽然说过—\par

\Quote{灵魂消逝如幻影，\newline{}}\par

在另一处却说：—\par

\Quote{无体的灵魂进入阴府；\newline{}}\par

又说：—\par

\Quote{使我速过阴府门，\newline{}将我埋葬。\newline{}}\par

至于你所读的其他作家，我想你已足够准确地知道他们如何表达。但这一切，凡寻求天主的智慧，并通过信德、正义和善行蒙天主喜悦的人都会明白。因为我们之前提到的先知之一，名叫\Person{欧瑟亚}，说：\BibleQuote{谁明智，必明了这些事；谁聪敏，必认识它们！因为主的道路是正直的，义人在其中行走，但背逆者必在其中跌倒。}\BibleRef{欧 14:9}因此，愿意学习的人应当多学。你要努力更常与我相见，好通过亲耳听讲，准确查明真理。\par

\SourceRecord{Translated by Marcus Dods.；Ante-Nicene Fathers；Vol. 2.；Edited by Alexander Roberts, James Donaldson, and A. Cleveland Coxe.；Buffalo, NY: Christian Literature Publishing Co.,；1885.；<http://www.newadvent.org/fathers/02042.htm>.}

% structure: p0001 source=h1 target=section reason=page-title

\section{致\Person{奥托里科}，第三卷}

% structure: p0002 source=h2 target=subsection reason=relative-heading-rank; base=h2

\subsection{第一章 \Person{奥托里科}尚未信服}

\Person{德奥斐洛}致候\Person{奥托里科}。鉴于著作家们喜爱为虚荣编写众多书籍——有些涉及神祇、战争和编年史，有些则涉及无益的传说和其他虚妄的劳作，你至今也惯常投身其中，并不吝于忍受那种辛劳；但尽管你曾与我交谈，仍认为真理之言是无聊的故事，并以为我们的著作是近代新作——因此，我也将不辞辛劳，在\OriginalTerm{天主}{God}的帮助下，向你简明陈述我们书籍的古老性，以寥寥数语提醒你，使你不吝于阅读，而能认识其他作者的愚妄。\par

% structure: p0004 source=h2 target=subsection reason=relative-heading-rank; base=h2

\subsection{第二章 世俗作者无法认识真理}

因为著作者本当亲自目睹其所声称之事，或从目击者那里准确查证；因为撰写未经证实之事的人只是白费力气。\Person{荷马}编造特洛伊战争、欺骗众人，有何益处？\Person{赫西俄德}编写他所称神明的谱系，有何益处？\Person{俄耳甫斯}的三百六十五位神明——临终时却又否认，在其训言中主张唯有\OriginalTerm{一位天主}{one God}——有何益处？\Person{阿拉托斯}的宇宙圈球面图，或与他持同样学说的人，除了人间的荣耀，有何益处？连这荣耀他们也未得应得的。他们说了什么真理？\Person{欧里庇得斯}、\Person{索福克勒斯}或其他悲剧作家的悲剧有何益处？\Person{米南德}、\Person{阿里斯托芬}或其他喜剧作家的喜剧呢？\Person{希罗多德}、\Person{修昔底德}的历史呢？\Person{毕达哥拉斯}的赫拉克勒斯神殿和石柱呢？\Person{狄奥基尼}的犬儒哲学呢？\Person{伊壁鸠鲁}主张没有天主眷顾，有何益处？\Person{恩培多克勒}教导无神论呢？\Person{苏格拉底}以狗、鹅、悬铃木、被雷电击中的\Person{阿斯克勒庇俄斯}以及他所呼求的精灵起誓，有何益处？他为何情愿赴死？他期望死后得到何种奖赏？\Person{柏拉图}的教养体系对他有何益处？\Emph{其他}哲学家从他们的学说中得到什么益处？——我不必一一列举，因为人数众多。但我们说这些，是为了显明他们无益且不敬天主的见解。\par

% structure: p0006 source=h2 target=subsection reason=relative-heading-rank; base=h2

\subsection{第三章 他们的矛盾}

因所有这些人都迷恋虚浮的空名，既不自知真理，也未引导他人认识真理；他们自己的话语就证明他们自相矛盾，且大多数人推翻了自身的学说。因为他们不仅相互驳斥，有些人甚至否定了自己的教导，以致其名声终结于羞耻和愚妄，为明达之人所谴责。因为他们或论及神明，随后又教导无神论；或谈论世界的创造，最后又说万物是自发产生的。甚至当谈论天主眷顾时，却又教导世界不受眷顾。不仅如此，当他们尝试写作关于高尚行为时，却教导行淫乱、奸淫和通奸，并引入可憎的、不可言说的邪恶。他们宣称他们的神明率先犯下不可言说的奸淫和骇人听闻的宴饮。谁不歌唱\Person{萨图尔努斯}吞噬自己的儿女，\Person{朱庇特}吞下其子\Person{墨提斯}，并为众神准备可怕的宴席——据说瘸腿铁匠\Person{伏尔甘}在那里侍候——以及\Person{朱庇特}不仅娶了自己的妹妹\Person{朱诺}，还用污秽的口行可憎的恶事？至于他的其余事迹，诗人所歌唱的许多，想必你都知道。我何须再叙述\Person{涅普顿}、\Person{阿波罗}、\Person{巴克斯}、\Person{赫拉克勒斯}、爱胸部的\Person{密涅瓦}和无耻的\Person{维纳斯}的事迹呢？因为在别处我们已经更准确地讲述了这些。\par

% structure: p0008 source=h2 target=subsection reason=relative-heading-rank; base=h2

\subsection{第四章 \Person{奥托里科}如何被针对基督徒的虚假指控误导}

我本无需驳斥这些，只是见你仍对真理之言存疑。你虽明智，却甘愿忍受愚昧人。否则，你不会被无头脑之人打动，屈服于空言，相信那普遍流传的谣言——不敬\OriginalTerm{天主}{God}的口唇以此诬蔑我们这些崇拜\OriginalTerm{天主}{God}、称为基督徒的人，声称我们所有人的妻子被共有并随意使用，我们甚至与自己的姐妹乱伦，最不虔敬、最野蛮的是，我们吃人肉。此外，他们还说我等的教义只是最近才出现，我们没有任何证据证明我们所接受的真理或我们的教导，反而我们的教义是愚拙。那么，我尤其惊讶，你在其他事情上勤勉钻研，洞察一切，却在我们的事上掉以轻心。因为，若有可能，你必不惜彻夜泡在图书馆。\par

% structure: p0010 source=h2 target=subsection reason=relative-heading-rank; base=h2

\subsection{第五章 哲学家鼓吹同类相食}

既然如此，你读了许多书，对于\Person{芝诺}、\Person{第欧根尼}和\Person{克莱安西斯}的训诲，即他们的著作所包含的，强令人吃人肉：父亲被其亲生儿女煮熟吃掉；若有人拒绝或不接受这无耻食物的一部分，那不愿吃的人自己反要被吃，你有什么看法？还有比这更不敬神的话——即\Person{第欧根尼}的话，他教导儿童把自己的父母献为祭品并吃掉。历史学家\Person{希罗多德}不是也叙述了\Person{冈比西斯}在杀了\Person{哈尔帕戈斯}的儿女之后，将他们煮熟，摆在他们父亲面前作为餐食吗？而且，他还叙述说，在印度人中，父母被自己的儿女吃掉。哦！那些记录，甚至灌输这类事的人的渎神教导！哦！他们的邪恶与不敬！哦！那些如此精确地探究哲学并自诩为哲学家的人的思想！因为教导这些学说的人已使世界充满不义。\par

% structure: p0012 source=h2 target=subsection reason=relative-heading-rank; base=h2

\subsection{第六章　哲学家的其他观点。}

至于不法行为，那些盲目闯入哲学歌队的人，几乎众口一词。当然，\Person{柏拉图}——先提及他，因他似乎是他们中最受尊敬的哲学家——在他的第一部著作\WorkTitle{理想国}中，直接地，似乎立法规定所有人共妻，借用朱庇特之子与克里特立法者的先例，以便在此借口下，最佳之人能产生众多后代，而疲惫不堪的人能通过此类交合得到安慰。\Person{伊壁鸠鲁}本人，在教导无神论的同时，也教导与母亲和姐妹乱伦，并且这是违反禁止此类行为的法律；因为\Person{梭伦}明确就此立法，规定已婚父母所生子女为合法，不得生于奸淫，以使无人将非其父者尊为父，或因其不知而羞辱真正的生父。罗马人和希腊人的其他法律也禁止这些事。那么，为什么\Person{伊壁鸠鲁}和斯多葛派教导乱伦和鸡奸，并用这些学说塞满图书馆，以致从幼年就学习这种不法交合？我还何必花费更多时间论他们，既然连他们称为神的那类事物他们也讲述类似之事？\par

% structure: p0014 source=h2 target=subsection reason=relative-heading-rank; base=h2

\subsection{第七章　关于诸神的不同学说。}

因为他们宣称这些是神之后，又使他们毫无价值。有人说他们由原子构成；又有人说他们最终归于原子；他们说神并不比人更有能力。\Person{柏拉图}，虽他说这些是神，却认为他们由物质构成。\Person{毕达哥拉斯}，在如此辛劳地研究神，并四处奔波（寻求信息）之后，最终断定万物皆自然自发生成，神不关心人类。阿卡德米学派的\Person{克利托马库斯}引入了多少无神论观点，（我无需赘述）。\Place{阿布德拉}的\Person{克里提亚斯}和\Person{普罗泰戈拉}不是说过：\Quote{因为我既不能肯定诸神存在，也不能解释他们的本性；有许多事阻碍我}吗？至于最无神论者\Person{欧赫墨罗斯}的观点，更是多余。因为他虽对神作了许多大胆断言，最终却完全否认神的存在，认为万物由自发运动支配。而\Person{柏拉图}，他曾大谈神之合一及人的灵魂，断言灵魂不朽，他自己后来岂不又前后矛盾地主张有些灵魂进入其他人身，而有些则进入非理性动物？他的学说怎能不显得可怕而怪诞——至少对稍有判断力的人而言——一个人从前是人，后来却成为狼、狗、驴或其他非理性畜生？\Person{毕达哥拉斯}也发现他散布类似的胡言，此外还否定了眷顾。那么我们该相信谁？喜剧诗人\Person{菲利蒙}说——\par

\Quote{赞美并侍奉诸神者拥有美好的希望；}\par

抑或是我们已提及的那些人——\Person{欧赫墨罗斯}、\Person{伊壁鸠鲁}、\Person{毕达哥拉斯}，以及其他否认诸神应受敬拜并废除眷顾的人？关于神和眷顾，\Person{阿里斯顿}说：——\par

\Quote{鼓起勇气：神必保守\newline{} 并大大帮助所有配得之人。\newline{} 若忠诚之人不得升迁，\newline{} 试问行善有何益处？\newline{} 诚然——我常见义人\newline{} 遭遇不幸，被一切尊荣抛弃；\newline{} 而那些以自我晋升为目标的人，\newline{} 此生常已获得所求。\newline{} 但我们须看得更远，等待结局，\newline{} 万物所趋向的终局。\newline{} 并非如虚妄恶者所言，\newline{} 我们被失控的命运牵引；\newline{} 并非盲目的机运统治世界，\newline{} 也不是万物不受控地向前冲。\newline{} 恶人以这信念自欺；\newline{} 但要确信，所有奖赏中最大的\newline{} 仍为圣洁生活者存留；\newline{} 而眷顾将给予恶人\newline{} 他们所应得的公正报应，\newline{} 以及各恶行当受的惩罚。}\par

人们可以看到，其他人——实际上几乎是多数人——关于神和眷顾所说的，彼此多么不一致。因为有些人完全取消了神和眷顾；而另一些人则肯定神，并宣称万事皆由眷顾支配。因此，聪慧的听者和读者必须仔细留意他们的表述；正如\Person{西米洛斯}所说：\Quote{诗人的习惯是用一个共同名称称呼极恶者和卓越者；因此我们必须辨别。} 又如\Person{菲利蒙}所说：\Quote{一个坐着只听而不思考的愚昧人是个麻烦的特征；因为他不会自责，他是如此愚蠢。} 因此我们必须留意，并思考所说的话，批判性地探究哲学家和诗人所表述的内容。\par

% structure: p0020 source=h2 target=subsection reason=relative-heading-rank; base=h2

\subsection{第八章　异教作家归于诸神的邪恶。}

他们既否认有神，又承认神存在，并且说自己曾犯下极恶的行径。首先，诗人们动听地歌颂朱庇特（\Person{Jove}）的恶行。而说了大量废话的克吕西波（\Person{Chrysippus}），难道不被发现公开说朱诺（\Person{Juno}）与朱庇特（\Person{Jupiter}）有最污秽的交合吗？我何必一一列举所谓的众神之母的污秽，或是嗜饮人血的拉提亚里斯朱庇特（\Person{Jupiter Latiaris}），或是被阉割的阿提斯（\Person{Attis}）；或是那位绰号悲剧作家的朱庇特（\Person{Jupiter}），他如何玷污自己（如他们所说），如今又在罗马人中受敬拜为神？我不提安提诺乌斯（\Person{Antinous}）的神庙，以及你们称为神的其他人。因为讲给明理的人听，只会引起嘲笑。那些炮制哲学，否认神的存在，或主张混乱交合与兽性同居的人，他们自己的教训正定了自己的罪。此外，我们从他们撰写的著作中发现，吃人肉在他们当中被接受；他们还记载，他们所尊为神的人首先做了这些事。\par

% structure: p0022 source=h2 target=subsection reason=relative-heading-rank; base=h2

\subsection{第九章 基督教关于天主及其法律之训导}

现在我们也承认天主存在，但祂是独一之天主，是这宇宙的创造者、造作者和塑造者；我们知道万物都由祂的眷顾安排，但唯独由祂安排。我们领受了神圣的法律；而真正的天主就是立法者，祂教导我们行义、虔诚和行善。关于虔诚，祂说：\BibleQuote{你不可有别的神在我面前。不可为你制造任何雕像，或上天下地、地上或水中的任何肖像；不可俯伏朝拜它们，也不可事奉它们，因为我上主你的天主。}\BibleRef{出 20:3}关于行善，祂说：\BibleQuote{当孝敬你的父亲和母亲，使你在上主你的天主赐给你的地上得享幸福，并延长你的寿命。}关于正义，祂说：\BibleQuote{不可奸淫。不可杀人。不可偷盗。不可作假见证陷害你的邻人。不可贪恋你邻人的妻子，不可贪恋你邻人的房屋、田地、仆婢、牛驴、牲畜，或属于你邻人的任何东西。不可在穷人的诉讼中屈枉正直。从一切不义之事要远离。不可杀害无辜和正直的人；不可使恶人得称为义；也不可收受贿赂，因为贿赂能使明眼人失明，曲解正直的话。}\BibleRef{出 23:6}这神圣的法律，梅瑟（\Person{Moses}）——他也是天主的仆人——被立为传报者，传报给全世界，尤其是给希伯来人，他们也被称为犹太人，埃及王古时曾奴役他们，他们是亚巴郎（\Person{Abraham}）、依撒格（\Person{Isaac}）和雅各伯（\Person{Jacob}）这些虔诚圣洁之人的正义苗裔。天主记念他们，藉梅瑟的手行了奇妙惊人的奇迹，拯救他们，领他们出埃及，带领他们经过所谓的旷野；又使他们定居在客纳罕（\Place{Canaan}）地，后来称为犹太（\Place{Judaea}），赐给他们法律，教导他们这些事。这伟大而奇妙的、趋向一切正义的法律，其十项纲要正如我们刚才所复述的。\par

% structure: p0024 source=h2 target=subsection reason=relative-heading-rank; base=h2

\subsection{第十章 善待外方人}

因此，他们曾在埃及地作外方人，本是生于加色丁（\Place{Chaldaea}）的希伯来人——因为那时发生饥荒，他们被迫移居埃及去购粮，并在那里寄居一时；这些事按天主的预言临到他们——他们在埃及寄居了四百三十年，当梅瑟（\Person{Moses}）要领他们出埃及进入旷野时，天主藉法律教训他们说：\BibleQuote{不可苛待外方人，因为你们知道外方人的心情；因为你们自己曾在埃及地作过外方人。}\BibleRef{出 22:21}\par

% structure: p0026 source=h2 target=subsection reason=relative-heading-rank; base=h2

\subsection{第十一章 论悔改}

当人民违犯天主赐给他们的法律时，天主良善怜悯，不愿毁灭他们。除了赐给他们法律外，后来也派遣先知到他们中间，从他们的弟兄中兴起，教训和提醒他们法律的内容，引导他们悔改，使他们不再犯罪。如果他们坚持作恶，祂就预先警告他们，他们必被交与世上万国受奴役；这事已经应验在他们身上，是显而易见的。关于悔改，依撒意亚（\Person{Isaiah}）先知对众人，特别是对那人民说：\BibleQuote{趁上主可找到的时候，你们应寻求祂；趁祂还在近处的时候，你们应呼求祂。恶人应离开自己的行径，不义的人应抛弃自己的思念，归向上主，好得到仁慈，因为祂必广施赦免。}\BibleRef{依 55:6}另一位先知厄则克耳（\Person{Ezekiel}）说：\BibleQuote{若恶人从他所作的一切罪恶上转离，谨守我的一切法令，秉公行义，他必生存，不至死亡。他所行的一切不义，不再被记念；他因所行的正义，必得生存。因为我不喜欢恶人丧亡，却喜欢恶人离开他的罪行而得以生存。}\BibleRef{则 18:21}依撒意亚又说：\BibleQuote{你们这些深谋远虑行恶的人，回转吧，好得救。}\BibleRef{依 31:6}另一位先知耶肋米亚（\Person{Jeremiah}）说：\BibleQuote{归向上主你们的天主，像摘葡萄的人回到篮筐，你们必会找到仁慈。}\BibleRef{耶 6:9}因此，圣经中关于悔改的言论很多，甚至无数，因为天主常愿人类转离一切罪恶。\par

% structure: p0028 source=h2 target=subsection reason=relative-heading-rank; base=h2

\subsection{第十二章 论正义}

况且，关于法律所吩咐的义德，先知与福音中都有一致的论述，因为他们都由同一的天主圣神所感动而发言。依撒意亚因此这样说：\BibleQuote{从你们的心灵中除去你们行为的邪恶；学习行善，寻求正义，压迫者，为孤儿伸冤，为寡妇辩护。}\BibleRef{依 1:16} 同一先知又说：\BibleQuote{解开一切邪恶的绳索，废除一切压迫的契约，让受压迫者自由离去，撕毁一切不义的文契。把你的食物分给饥饿的人，将无家可归的穷人领到你家中。见到赤身露体的人，要给他遮盖，不要躲避你的骨肉。那时，你的光明要如晨光迸发，你的健康要迅速滋长，你的正义要在你面前行走。}\BibleRef{依 58:6} 同样，耶肋米亚也说：\BibleQuote{站在路上，察看并询问哪一条是上主你们的天主的美好的道路，然后行走其上，你们的心灵必得安息。施行公正的审判，因为这是上主你们的天主的旨意。}\BibleRef{耶 6:16} 欧瑟亚也这样说：\BibleQuote{持守正义，亲近你们的天主，他奠定了诸天，创造了大地。}\BibleRef{欧 12:6} 另一位先知岳厄尔与这些话一致地说：\BibleQuote{聚集民众，圣化集会，召集长老，聚集怀抱的孩童；让新郎走出他的洞房，新娘走出她的内室，并恳切地向主你们的天主祈祷，使他怜悯你们，涂抹你们的罪过。}\BibleRef{岳 2:16} 同样，另一位先知匝加利亚也说：\BibleQuote{万军的上主这样说：你们要执行真正的审判，人与人之间要显示仁慈和怜悯；不要欺压寡妇、孤儿和外方人；你们心中不要彼此图谋恶事——万军的上主说。}\BibleRef{匝 7:9}\par

% structure: p0030 source=h2 target=subsection reason=relative-heading-rank; base=h2

\subsection{第十三章 论贞洁}

论到贞洁，圣言教导我们不仅不可在行为上犯罪，甚至不可在思想中犯罪，不可在心里存有任何恶念，也不可用眼目贪恋他人的妻子。为此，身为君王和先知的撒罗满说：\BibleQuote{你的眼睛要向前直视，你的眼睑要直对前方；为你的脚修平道路。}\BibleRef{箴 4:25} 福音之声更严厉地教导贞洁，说：\BibleQuote{凡注视不是他自己妻子的女人而贪恋她的，已经在心里与她犯了奸淫。}\BibleRef{玛 5:28} 又说：\BibleQuote{凡娶那被丈夫休弃的妇人，也是犯奸淫；凡休妻的，除非为淫乱的缘故，就是使她犯奸淫。}\BibleRef{玛 5:32} 因为撒罗满说：\BibleQuote{人若怀里揣火，衣服岂能不烧呢？人若在热炭上行走，脚岂能不灼伤呢？同样，与有夫之妇亲近的人，必不能无罪。}\BibleRef{箴 6:27}\par

% structure: p0032 source=h2 target=subsection reason=relative-heading-rank; base=h2

\subsection{第十四章 论爱仇敌}

而且，我们不仅应对同族的人心怀善意（正如某些人所想的），依撒意亚先知说：\BibleQuote{对那些恨你们并把你们赶出去的人说：你们是我们的弟兄，使上主的名得荣耀，并在他们的喜乐中显明出来。}\BibleRef{依 66:5} 福音说：\BibleQuote{爱你们的仇敌，为那迫害你们的祈祷。因为如果你们只爱那爱你们的人，有什么赏报呢？就是强盗和税吏也这样做。}\BibleRef{玛 5:44} 行善的人也不可自夸，免得成为取悦于人者。因为经上说：\BibleQuote{不要让你的左手知道右手所行的。}\BibleRef{玛 6:3} 此外，关于服从掌权者和执政者并为他们祈祷，神圣的话给我们指示，为使我们\BibleQuote{能平安宁静地度日}\BibleRef{弟前 2:2}。并教导我们对众人当尽本分：\BibleQuote{该恭敬的，恭敬他；该畏惧的，畏惧他；该纳粮的，纳粮；不欠任何人什么，但要爱众人。}\BibleRef{罗 13:7}\par

% structure: p0034 source=h2 target=subsection reason=relative-heading-rank; base=h2

\subsection{第十五章 基督徒的无辜辩护}

因此，请想一想：那些教导这些道理的人，怎么可能过着漫不经心的生活，参与非法的苟合，甚至（最邪恶不过）吃人肉呢？尤其我们连观看角斗士表演都被禁止，以免我们成为杀人的同谋和帮凶。我们也不可观看其他表演，免得我们的眼目和耳朵被污染，参与其中所歌颂的话语。因为如果有人谈论食人肉，在这些表演中，有提厄斯忒斯和忒柔斯的儿女被吃；至于奸淫，无论人间还是神明（人们用优雅的语言歌颂他们，以赢得荣誉和奖品），都成为戏剧的主题。但基督徒绝不会做出如此行为；因为他们身上居住着节制，实行自律，遵守一夫一妻制，守护贞洁，根除不义，铲除罪恶，施行正义，奉行法律，从事敬拜，承认天主；真理引导，恩宠护卫，平安遮盖他们；圣言指引，智慧教导，生活导向，天主为王。因此，尽管关于我们的生活方式和创造万有的天主的律例，我们还有许多话可说，但我们认为目前所提醒你们的已足以促使你们研究这些事，尤其是你们现在可以自行阅读〔我们的著作〕，正如你们一向喜好求知，在这方面的学习也会持续不断。\par

% structure: p0036 source=h2 target=subsection reason=relative-heading-rank; base=h2

\subsection{第十六章 哲学家们不确定的猜测}

但我现在愿意，在天主的帮助下，给你们一个关于历史时期的更精确的证明，使你们看到我们的教义并非现代或虚构的，而是比所有在不确定中写作的诗人和作者更古老、更真实。因为有些人坚持世界是不受造的，便陷入了无限；另一些人断言世界是被造的，就说已经过去了153,075年。这是埃及人阿波罗尼奥所说的。而被认为是最有智慧的希腊人柏拉图，陷入了何等的胡言乱语？因为在他的题为 \WorkTitle{理想国} 的书中，我们明确地发现他说：\Quote{如果事物在所有时期中都保持着现在的安排，那么什么时候能发现任何新事物呢？因为一万万年的时间没有记录就过去了，而从一些事物被代达罗斯、奥菲斯和帕拉墨得斯发现以来，也已经过去了一千年或两千年。}而当他说这些事发生时，他暗示从洪水到达达罗斯过去了一万年。在他关于世界的城市、定居点和民族说了很多之后，他承认这些是猜测性的。因为他说：\Quote{那么，我的朋友，如果某个神应许我们，如果我们试图对立法进行考察，现在所说的这些……}等等，这表明他是凭猜测说话的；如果是凭猜测，那么他所说的就不真实。\par

% structure: p0038 source=h2 target=subsection reason=relative-heading-rank; base=h2

\subsection{第十七章. 基督徒的精确信息。}

因此，他更应该在这立法之事上成为天主的门徒，正如他亲自承认的那样，除非天主通过法律教导他，否则他无法获得准确的信息。诗人荷马、赫西俄德和奥菲斯不也宣称他们自己曾受到天主眷顾的教导吗？此外，据说在你们的作者中有先知和预言家，那些受他们教导的人写得准确。那么，我们这些受圣先知教导的人，他们被天主圣神所充满，岂不更应知道真理！为此，所有先知和谐一致地说话，彼此同意，并预言了将在全世界发生的事。因为预言和已经完成的事件本身的发生，应该向那些喜爱信息，更热爱真理的人证明，他们关于洪水之前时代和时期的宣告确实是真实的：即世界被创造以来至今的年代如何运行，以显示你们作者的荒谬谎言，并表明他们的陈述不真实。\par

% structure: p0040 source=h2 target=subsection reason=relative-heading-rank; base=h2

\subsection{第十八章. 希腊人关于洪水的错误。}

因为正如我们上面所说，柏拉图在证明发生过洪水时，说它没有覆盖整个大地，只覆盖了平原，而那些逃到最高山上的人得救了。但其他人说存在丢卡利翁和皮拉，他们被保存在一个箱子里；丢卡利翁从箱子里出来后，把石头扔在身后，石头产生了人；因此他们称大众为\Quote{人民}。另一些人又说，在第二次洪水中存在克莱墨诺斯。从已经说过的来看，显然那些写这类东西并如此无益地从事哲学的人是可怜的，非常不敬和愚昧的人。但我们的先知、天主的仆人梅瑟，在描述世界的起源时，讲述了洪水是如何降临到地上的，还告诉了我们洪水的细节是如何发生的，没有讲皮拉、丢卡利翁或克莱墨诺斯的寓言；当然也没有说只有平原被淹没，只有那些逃到山上的人得救。\par

% structure: p0042 source=h2 target=subsection reason=relative-heading-rank; base=h2

\subsection{第十九章. 洪水的准确记载。}

他也没有说有过第二次洪水；相反，他说水灾将不再临到世界；确实没有，也永远不会再有。他说有八个人在方舟里得救，这方舟是根据天主的指示准备的，不是由丢卡利翁，而是由诺厄建造的；这个希伯来词意思是\Quote{安息}，正如我们在别处指出的那样，诺厄在向当时活着的人宣布洪水即将来临时，向他们预言说：\Quote{来吧，天主呼召你们悔改。}为此，他恰当地被称为丢卡利翁。这个诺厄有三个儿子（正如我们在第二卷中提到的），名叫闪、含和雅弗；他们有三个妻子，每人一个妻子；每人及其妻子。这个人有些人称之为欧努科斯。因此，方舟里所有的八个人都得救了。梅瑟表明洪水持续了四十天四十夜，暴雨从天倾泻，深渊的泉源也裂开，水涨过所有高山十五肘。因此，当时所有人类的种族都被毁灭了，只有那些在方舟里受保护的得救了；我们已经说过，他们是八个人。方舟的残骸至今在阿拉伯山脉中仍可见到。以上就是洪水历史的概要。\par

% structure: p0044 source=h2 target=subsection reason=relative-heading-rank; base=h2

\subsection{第二十章. 梅瑟的古老性。}

\Person{梅瑟}成了犹太人的领袖，如我们已说过的那样，被法郎王\Person{阿玛西斯}从埃及地驱逐出去；据说\Person{阿玛西斯}在驱逐百姓后统治了二十五年又四个月，按\Person{玛奈托}所假定。他之后是\Person{赫布隆}，十三年。他之后是\Person{阿美诺菲斯}，二十年又七个月。他之后是其妹\Person{阿梅萨}，二十一年又一个月。她之后是\Person{梅弗勒斯}，十二年又九个月。他之后是\Person{梅特拉穆托西斯}，二十年又十个月。他之后是\Person{提特莫西斯}，九年又八个月。他之后是\Person{达姆菲诺菲斯}，三十年又十个月。他之后是\Person{奥鲁斯}，三十五年又五个月。他之后是其女，十年又三个月。她之后是\Person{梅尔克勒斯}，十二年又三个月。他之后是其子\Person{阿尔迈斯}，三十年又一个月。他之后是\Person{米亚姆穆斯}之子\Person{梅塞斯}，六年又两个月。他之后是\Person{拉美西斯}，一年又四个月。他之后是\Person{阿美诺菲斯}，十九年又六个月。他之后是其子\Person{托厄苏斯}和\Person{拉美西斯}，十年，据说他们拥有一支庞大的骑兵队和海军装备。希伯来人，确实，按他们自己的独特历史，那时已移居埃及地，并被王\Person{特特莫西斯}奴役，如上所述，为他建造了坚固的城市：\Place{培通}、\Place{拉美西斯}和\Place{昂}，即\Place{赫利奥波利斯}；因此，希伯来人，也是我们的祖先，并从他们那里我们得到了那些神圣的经书，比所有作者都古老，如已说过的那样，被证明比当时在埃及人中著名的城市更古老。这地因王\Person{塞托斯}而被称为埃及。因为据说\Person{塞托斯}这个词发音为\Quote{埃及}。\Person{塞托斯}有一个兄弟，名叫\Person{阿尔迈斯}。他被称为\Person{达那俄斯}，正是那位从\Place{埃及}前往\Place{阿尔戈斯}的人，其他作者提及他时认为他非常古老。\par

% structure: p0046 source=h2 target=subsection reason=relative-heading-rank; base=h2

\subsection{第二十一章　论\Person{玛奈托}之不准确}

在\Place{埃及}人中散播了大量无稽之谈、甚至不敬地指控\Person{梅瑟}和与他同行的希伯来人因癞病被逐出\Place{埃及}的\Person{玛奈托}，未能给出准确的年代陈述。因为当他说他们是牧羊人、是埃及人的敌人时，他确实说了真话，因为他被迫如此。因为我们寄居\Place{埃及}的祖先确实是牧羊人，而非癞病人。因为他们进入称为\Place{耶路撒冷}的地方，后来也住在那里，众所周知，他们的司铎按照天主的任命，继续在圣殿里，并在那里医治各种疾病，以致他们治愈了癞病人和一切不洁。圣殿是由犹太王\Person{撒罗满}建造的。从\Person{玛奈托}自己的陈述来看，他的年代错误是明显的（关于驱逐他们的王、名叫\Person{法郎}，也是如此。因为他不再统治他们。因为他追赶希伯来人，他和他的军队被淹没在\Place{红海}中。他还有更进一步的错误，说牧羊人曾与埃及人作战）。因为他们离开了\Place{埃及}，此后住在现今称为\Place{犹太}的地方，比\Person{达那俄斯}到达\Place{阿尔戈斯}早三百一十三年。大多数人认为他比任何其他希腊人都更古老，这是显而易见的。因此，\Person{玛奈托}通过他自己的著作，不情愿地向我们宣告了两件真实的事情：第一，承认他们是牧羊人；第二，说他们离开了\Place{埃及}地。因此，即使从这些著作中，\Person{梅瑟}及其追随者被证明比\Place{特洛伊}战争早九百年甚至一千年。\par

% structure: p0048 source=h2 target=subsection reason=relative-heading-rank; base=h2

\subsection{第二十二章　圣殿之古老}

然后关于\Place{犹太}地圣殿的建造，这是\Person{撒罗满}王在犹太人出\Place{埃及}后五百六十六年建造的，在\Place{提洛}人中有一份关于圣殿如何建造的记录；在他们的档案中保存着一些文件，其中证明圣殿在\Place{提洛}人建立\Place{迦太基}之前一百四十三年零八个月就已存在（这份记录是由\Person{希兰}——这是\Place{提洛}王的名字，\Person{阿比玛鲁斯}之子——所作，因为\Person{希兰}和\Person{撒罗满}之间存在世袭的友谊，同时因为\Person{撒罗满}拥有卓越的智慧。因为他们不断互相讨论难题。这在他们至今仍保存在\Place{提洛}人中的通信以及他们之间往来的文件中得到了证明）；正如\Place{以弗所}的\Person{墨涅安得耳}在叙述\Place{提洛}王国的历史时所记录的那样，他这样说道：\Quote{当\Place{提洛}王\Person{阿比玛鲁斯}死后，其子\Person{希兰}继位。他活了五十三年。\Person{巴佐鲁斯}继他之后，活了四十三年，统治了十七年。之后是\Person{梅图阿斯塔图斯}，活了五十四年，统治了十二年。之后是其弟\Person{阿塔里穆斯}，活了五十八年，统治了九年。他被其兄弟\Person{赫勒斯}所杀，\Person{赫勒斯}活了五十年，统治了八个月。他被\Person{阿斯塔特}神的司铎\Person{犹托巴鲁斯}所杀，\Person{犹托巴鲁斯}活了四十年，统治了十二年。其子\Person{巴佐鲁斯}继位，活了四十五年，统治了七年。其子\Person{梅滕}继位，活了三十二年，统治了二十九年。\Person{皮格马利翁}，\Person{皮格马利乌斯}之子，继位，活了五十六年，统治了七年。在他统治的第七年，其妹逃往\Place{利比亚}，建造了那座至今仍称为\Place{迦太基}的城市。}因此，从\Person{希兰}统治到\Place{迦太基}建立的全部时期总计一百五十五年零八个月。而在\Person{希兰}统治的第十二年，\Place{耶路撒冷}的圣殿被建造。所以从建造圣殿到\Place{迦太基}建立的全部时间是一百四十三年零八个月。\par

% structure: p0050 source=h2 target=subsection reason=relative-heading-rank; base=h2

\subsection{第二十三章　先知较\Place{希腊}作家更古老}

所以，以上有关腓尼基人和埃及人的见证，以及埃及人\Person{玛内托}和\Place{厄弗所}人\Person{默南德尔}所写的年代记录，还有撰写犹太人反抗\Place{罗马}战争史的作者\Person{若瑟夫}的记述，就足以作为证据了。因为这些十分古老的记录证明，其他民族的著作比通过\Person{梅瑟}赐给我们的著作，甚至比后来的先知们更为晚近。因为最后一位被称为\Person{匝加利亚}的先知，与\Person{达理阿}王同时代。但就连立法者们本人，也都是在那个时期之后才立法的。假如有人提及\Place{雅典}人\Person{索伦}，他生活在\Person{居鲁士}王和\Person{达理阿}王时代，与首先提到的先知\Person{匝加利亚}同时期，而\Person{匝加利亚}在许多年前就是最后一位先知了。或者你提及立法者\Person{李库尔果}、\Person{德拉科}或\Person{米诺斯}，\Person{若瑟夫}在他的著作中告诉我们，圣书在古老性上先于他们，因为在\Person{朱庇特}统治\Place{克里特}人之前，甚至在\Place{特洛伊}战争之前，通过\Person{梅瑟}赐给我们的神圣法律就已经存在了。为了更精确地展示时代和日期，我们在天主的帮助下，现在不仅要叙述洪水之后的日期，还要叙述洪水之前的日期，以便尽可能算出全部年数的总数；一直追溯到天主之仆\Person{梅瑟}籍着圣神所记载的创造世界的开端。因为他首先论及世界的创造和起源，以及第一个人，然后按事件顺序叙述之后发生的一切，他也指明了洪水之前经过的年数。我祈求独一的天主恩佑，使我能够按照他的旨意准确地说出全部真理，以便你和每一位阅读此书的人都能因他的真理和恩宠而得到指引。那么，我将首先从记载的族谱开始，并以第一个人作为我叙述的开端。\par

% structure: p0052 source=h2 target=subsection reason=relative-heading-rank; base=h2

\subsection{第二十四章 自\Person{亚当}起的年代学。}

\Person{亚当}活了二百三十年，生子\Person{舍特}。\Person{舍特}二百零五年，生子\Person{厄诺士}。\Person{厄诺士}一百九十年，生子\Person{刻南}。\Person{刻南}一百七十年，生子\Person{玛拉肋耳}。\Person{玛拉肋耳}一百六十五年，生子\Person{雅勒得}。\Person{雅勒得}一百六十二年，生子\Person{哈诺客}。\Person{哈诺客}一百六十五年，生子\Person{默突舍拉}。\Person{默突舍拉}一百六十七年，生子\Person{拉默客}。\Person{拉默客}一百八十八年，生子\Person{诺厄}。\Person{诺厄}五百岁时生了\Person{闪}。\Person{诺厄}在世第六百年时，洪水来临。因此，直至洪水的总年数为二千二百四十二年。洪水之后，\Person{闪}一百岁生了\Person{阿帕革沙得}。\Person{阿帕革沙得}一百三十五岁生了\Person{撒拉}。\Person{撒拉}一百三十岁生子。其子\Person{厄贝尔}一百三十四岁。希伯来人由此得名。其子\Person{培肋格}一百三十岁生子。其子\Person{勒伍}一百三十二岁。其子\Person{色鲁格}一百三十岁。其子\Person{纳曷尔}七十五岁。其子\Person{特辣黑}七十岁。其子\Person{亚巴郎}，我们的圣祖，一百岁时生了\Person{依撒格}。因此，直至\Person{亚巴郎}，共有三千二百七十八年。上述的\Person{依撒格}六十岁时生子\Person{雅各伯}。\Person{雅各伯}直至移居\Place{埃及}，共一百三十年，我们已在上文提及。希伯来人在\Place{埃及}寄居了四百三十年；离开\Place{埃及}后，他们在所谓的旷野中度过了四十年。所有这些年数共计三千九百三十八年。那时\Person{梅瑟}去世，\Person{农}的儿子\Person{若苏厄}接替他的统治，治理了二十七年。\Person{若苏厄}之后，人民违犯了天主的诫命，被\Place{美索不达米亚}王\Person{哥托诺耳}统治了八年。然后，人民悔改，有了民长：\Person{敖特尼耳}四十年；\Person{厄革隆}十八年；\Person{厄胡得}八年。之后犯罪，被外族人压制二十年。然后\Person{德波辣}治理四十年。然后受\Place{米德扬}人统治七年。然后\Person{基德红}治理四十年；\Person{阿彼默肋客}三年；\Person{托拉}二十二年；\Person{雅依尔}二十二年。然后\Place{培肋舍特}人和\Place{阿孟}人统治他们十八年。此后\Person{依弗大}治理六年；\Person{厄斯邦}七年；\Person{厄隆}十年；\Person{阿贝冬}八年。然后外族人统治他们四十年。然后\Person{三松}治理二十年。然后他们之间有和平四十年。然后\Person{撒默辣}治理一年；\Person{厄里}二十年；\Person{撒慕尔}十二年。\par

% structure: p0054 source=h2 target=subsection reason=relative-heading-rank; base=h2

\subsection{第二十五章 自\Person{撒乌耳}至充军。}

在民长之后，他们有君王，第一位名叫\Person{撒乌耳}，执政二十年；然后是我们的先祖\Person{达味}，执政四十年。因此，从\Person{依撒格}到达味统治为止共有四百九十六年。在这些君王之后，\Person{撒罗满}执政，他也是奉天主的旨意，首先在\Place{耶路撒冷}建造圣殿的人；他执政四十年。在他之后是\Person{勒哈贝罕}，十七年；在他之后是\Person{阿彼雅}，七年；在他之后是\Person{阿撒}，四十一年；在他之后是\Person{约沙法特}，二十五年；在他之后是\Person{约兰}，八年；在他之后是\Person{阿哈齐雅}，一年；在他之后是\Person{阿塔里雅}，六年；在她之后是\Person{约史雅}，四十年；在他之后是\Person{阿玛责雅}，三十九年；在他之后是\Person{乌齐雅}，五十二年；在他之后是\Person{约堂}，十六年；在他之后是\Person{阿哈次}，十七年；在他之后是\Person{希则克雅}，二十九年；在他之后是\Person{默纳舍}，五十五年；在他之后是\Person{阿孟}，二年；在他之后是\Person{约史雅}，三十一年；在他之后是\Person{约阿哈次}，三个月；在他之后是\Person{约雅金}，十一年。然后是另一位\Person{约雅金}，三个月零十天；在他之后是\Person{漆德克雅}，十一年。在这些君王之后，百姓继续犯罪，不肯悔改，\Place{巴比伦}王，名叫\Person{拿步高}，按照\Person{耶肋米亚}的先知话，上来攻击\Place{犹太}。他将犹太人民迁到\Place{巴比伦}，并摧毁了\Person{撒罗满}所建的圣殿。在巴比伦充军期间，百姓度过了七十年。直到在\Place{巴比伦}寄居的时期，总共四千九百五十四年六个月零十天。正如天主曾藉\Person{耶肋米亚}先知预言百姓将被掳到\Place{巴比伦}，同样，祂也预先表明他们将在七十年后返回自己的土地。这七十年满了以后，\Person{居鲁士}成为\Place{波斯}王，他按照\Person{耶肋米亚}的先知话，在即位的第二年颁布谕令，命令他王国中的所有犹太人都返回本国，并为天主重建圣殿，就是前述\Place{巴比伦}王所拆毁的。此外，\Person{居鲁士}遵照天主的指示，命令他的卫兵\Person{萨贝撒尔}和\Person{米特利达特}，将\Person{拿步高}从\Place{犹太}圣殿中取走的器皿归还，并重新安放在圣殿中。因此，在\Person{达理阿}第二年，\Person{耶肋米亚}所预言的七十年得以应验。\par

% structure: p0056 source=h2 target=subsection reason=relative-heading-rank; base=h2

\subsection{第二十六章 希伯来著作与希腊著作的对比。}

由此可以看出，我们的圣经比希腊人、埃及人或任何其他历史学家的著作更古老、更真实。因为\Person{希罗多德}、\Person{修昔底德}、\Person{色诺芬}以及大多数其他历史学家，他们的叙述都是从\Person{居鲁士}和\Person{达理阿}统治时期左右开始的，无法准确讲述更早的古代时代。他们如果谈论\Person{达理阿}和\Person{居鲁士}这些蛮族国王，或者希腊的\Person{佐皮罗斯}和\Person{希庇亚斯}，或者\Place{雅典}人和\Place{拉栖代梦}人的战争，或者\Person{薛西斯}或\Person{保萨尼阿斯}的功绩——\Person{保萨尼阿斯}曾冒饿死在密涅瓦神庙中的危险——或者\Person{地米斯托克利}和\Place{伯罗奔尼撒}战争的历史，或者\Person{亚西比德}和\Person{特拉西布卢斯}的历史，那究竟揭示了什么重要的事情呢？因为我的目的不是提供空洞的言谈，而是要阐明从创世以来的年数，并谴责这些作者的虚妄劳作和琐碎，因为从洪水到现今既非像\Person{柏拉图}所说有两亿年，也非如我们已提及的埃及人\Person{阿波罗尼乌斯}所宣称的十五万五千六百二十五年；世界也不是非受造的，也不是像\Person{毕达哥拉斯}等人所梦想的那样万物自然发生；相反，世界受造而成，并由创造万有的天主的眷顾所统治；整个时间进程和年数对那些愿意服从真理的人都清晰可见。那么，以免我似乎只阐明到\Person{居鲁士}时代而忽略了随后的时期，好像无力展示它们，我将竭尽所能，依靠天主的助佑，对随后的时间进程加以说明。\par

% structure: p0058 source=h2 target=subsection reason=relative-heading-rank; base=h2

\subsection{第二十七章 罗马年代学至\Person{马尔库斯·奥勒留}之死。}

当时，居鲁士已统治了二十九年，在马萨格泰人之地被托米丽丝杀死，时在第六十二届\OriginalTerm{奥林匹亚纪}{Olympiad}；于是罗马人开始强大，因为天主加强他们。罗马由罗慕路斯建立，他是玛尔斯和伊利亚所传说的儿子，时在第七届\OriginalTerm{奥林匹亚纪}{Olympiad}，四月二十一日，当时一年按十个月计算。居鲁士既死，如我们所说，在第六十二届\OriginalTerm{奥林匹亚纪}{Olympiad}，这一年相当于\OriginalTerm{罗马建城纪元}{A.U.C.}220年；同年，塔奎尼乌斯·苏培布斯（与其绰号\Quote{\OriginalTerm{高傲}{Superbus}}）统治了罗马人。他是第一个放逐罗马人、败坏青年、使公民成为阉人，并且首先玷污处女然后将其嫁出的人。因此，他在罗马语中被称为\Quote{\OriginalTerm{高傲}{Superbus}}。因为他首先下令，凡向他致敬者，应由他人代受致敬。他统治了二十五年。此后，开始有年度执政官、保民官和营造官，共四百五十三年。列出这些官员的名字，我们认为冗长且多余。若有人急于了解，可从改名者克律塞鲁斯（\Person{Chryserus}，\OriginalTerm{命名者}{nomenclator}）编纂的表格中查得：他是奥雷利乌斯·维鲁斯（\Person{Aurelius Verus}）的被释奴，撰写了一部非常清晰的记录，涵盖从罗马建城到其恩主维鲁斯皇帝驾崩的所有事件，包括名字和日期。正如我们所说，年度长官统治了罗马人四百五十三年。此后，被称为皇帝的人按此顺序开始：首先，盖乌斯·尤利乌斯（\Person{Caius Julius}），统治三年四个月零六天；接着，奥古斯都（\Person{Augustus}），五十六年四个月零一天；提比略（\Person{Tiberius}），二十二年；接着另一位盖乌斯（\Person{Caius}），三年八个月零七天；克劳狄（\Person{Claudius}），二十三年八个月零二十四天；尼禄（\Person{Nero}），十三年六个月零五十八天；加尔巴（\Person{Galba}），二年七个月零六天；鄂图（\Person{Otho}），三个月零五天；维特里乌斯（\Person{Vitellius}），六个月零二十二天；韦斯巴芗（\Person{Vespasian}），九年十一个月零二十二天；提图斯（\Person{Titus}），二年零二十二天；图密善（\Person{Domitian}），十五年五个月零六天；涅尔瓦（\Person{Nerva}），一年四个月零十天；图拉真（\Person{Trajan}），十九年六个月零十六天；哈德良（\Person{Adrian}），二十年十个月零二十八天；安东尼（\Person{Antoninus}），二十二年七个月零六天；维鲁斯（\Person{Verus}），十九年零十天。因此，从凯撒到维鲁斯皇帝驾崩的时间为二百三十七年零五天。所以，从居鲁士之死以及塔奎尼乌斯·苏培布斯（\Person{Tarquinius Superbus}）的统治到维鲁斯皇帝驾崩，总年数为七百四十四年。\par

% structure: p0060 source=h2 target=subsection reason=relative-heading-rank; base=h2

\subsection{第二十八章 主要年代纪元}

至于从世界奠基以来的全部时间，就其主要纪元而言，可以如下追溯：从创造世界到洪水，共二千二百四十二年。从洪水到我们的先祖亚巴郎生了一个儿子，共一千零三十六年。从亚巴郎的儿子依撒格到百姓与梅瑟一同居住在旷野，共六百六十年。从梅瑟之死和农的儿子若苏厄的统治，到圣祖达味之死，共四百九十八年。从达味之死和撒罗满的统治，到百姓在巴比伦地寄居，共五百一十八年零六个月零十天。从居鲁士的治理到奥雷利乌斯·维鲁斯皇帝驾崩，共七百四十四年。从创造世界以来的所有年份，总数达五千六百九十八年，另加零月零日。\par

% structure: p0062 source=h2 target=subsection reason=relative-heading-rank; base=h2

\subsection{第二十九章 基督信仰的古老性}

这些时期以及上述所有事实综合起来，就能看到先知著作的古老性和我们教义的神性，即这教义不是新近的，我们的教条也不是某些人所认为的那样神话和虚假，而是非常古老和真实的。因为\OriginalTerm{塔鲁斯}{Thallus}提到了亚述人的王贝鲁斯和提坦的儿子萨图尔努斯，声称贝鲁斯与提坦人一同与朱庇特及其联盟中的所谓众神作战；此时他说古格斯战败，逃到了塔尔提苏斯。当时古格斯统治着那个地方，那时称为阿克特，现在称为阿提卡。至于其他国家和城市名称的来源，我们认为不必赘述，尤其是对您这样熟悉历史的人。梅瑟，不仅他，而且跟随他的大部分先知，都被证明早于所有作家，早于萨图尔努斯、贝鲁斯和特洛伊战争，这是显而易见的。因为根据塔鲁斯的历史，贝鲁斯被发现比特洛伊战争早三百二十二年。但我们在上面已表明，梅瑟大约比特洛伊陷落早九百年或一千年。由于萨图尔努斯和贝鲁斯活跃于同一时期，大多数人都不知道谁是萨图尔努斯，谁是贝鲁斯。有些人敬拜萨图尔努斯，称他为\OriginalTerm{贝尔}{Bel}或\OriginalTerm{巴尔}{Bal}，尤其东方国家的居民，因为他们不知道萨图尔努斯和贝鲁斯是谁。在罗马人中，他被称为萨图尔努斯，因为他们也不知道两者中哪个更古老——萨图尔努斯还是贝尔。至于奥林匹亚纪的开始，有些人说始于\OriginalTerm{伊菲托斯}{Iphitus}，但按其他人的说法，始于\OriginalTerm{利努斯}{Linus}，他亦被称为\OriginalTerm{伊利乌斯}{Ilius}。整个年份和奥林匹亚纪的顺序，我们已在上面表明。我认为，我已按我的能力，准确论述了你们习俗的无神性以及历史纪元的全部数字。因为即使我们在年代上犯了错误，比如五十年、一百年，甚至二百年，但也未像柏拉图、阿波罗尼乌斯和其他虚伪的作者迄今所写的那样，成千上万数万年。或许我们对全部年数的知识并非十分精确，因为神圣经书没有记载零月零日。但就我们所说的那些时期而言，我们有加色丁哲学家\OriginalTerm{贝罗苏斯}{Berosus}作证，他使希腊人了解了加色丁文献，并讲了一些关于洪水的事，还有许多其他历史要点，与梅瑟一致；而且与先知耶肋米亚和达尼尔，他也言辞一致。因为他提到在巴比伦人的王（他称之为\OriginalTerm{阿博巴索尔}{Abobassor}，希伯来人称之为\OriginalTerm{拿步高}{Nebuchadnezzar}）统治下犹太人发生的事。他还提到耶路撒冷圣殿；加色丁人的王如何使其荒凉，以及圣殿的根基在居鲁士统治第二年奠定，圣殿在达理阿统治第二年完工。\par

% structure: p0064 source=h2 target=subsection reason=relative-heading-rank; base=h2

\subsection{第三十章 为何希腊人未提及我们的历史}

但是希腊人并不提及那些讲述真理的历史：首先，因为他们自己只是在不久前才成为文字知识的分享者；他们自己也承认，声称文字是被发明的，有人说是在\Place{迦勒底}人当中，有人说在\Place{埃及}人当中，又有人说来自\Place{腓尼基}人。其次，因为他们犯了罪，并且仍在犯罪，不提及天主，反而提及虚妄无用的事。因为如此，他们最热诚地歌颂\Person{荷马}和\Person{赫西俄德}以及其他诗人，但对于那不朽的独一天主的荣耀，他们不仅避而不谈，甚至还亵渎；是的，他们迫害——并且每日仍在迫害——那些敬拜祂的人。不仅如此，他们甚至将奖品和荣誉颁给那些以和谐言词侮辱天主的人；但对于那些热心追求德性、实践圣洁生活的人，他们有的投石，有的处死，直到现今还对他们施以野蛮的酷刑。因此，这样的人必然失去了天主的智慧，并且没有找到真理。\par

如果你愿意，那么请仔细研读这些事情，好使你拥有真理的纲领和保证。\par

\SourceRecord{Translated by Marcus Dods.；Ante-Nicene Fathers；Vol. 2.；Edited by Alexander Roberts, James Donaldson, and A. Cleveland Coxe.；Buffalo, NY: Christian Literature Publishing Co.,；1885.；<http://www.newadvent.org/fathers/02043.htm>.}
